\chapter{机器人导航与具身智能}

\paragraph*{本章阅读锚点}
机器人场景中的 3DGS 不能只按“渲染好不好看”来读。对机器人而言，一个高斯场景要么作为真实世界的数字孪生，用于仿真、策略训练和评估；要么作为在线语义记忆，用于导航、探索和任务推理；要么作为可交互资产，与物理引擎、机械臂抓取、灵巧手或人形机器人动作连接。原始 3DGS 提供的是高保真外观和快速渲染，但机器人还需要尺度、碰撞几何、语义、可供性、物理参数、动作条件动态和闭环校正。本章因此按机器人用途展开：Real-to-Sim、导航记忆、机械臂操控、具身世界模型和主动建图。

\section{仿真与真实迁移}

Real-to-Sim 路线的锚点是“把真实场景变成可训练策略的仿真环境”。普通 3DGS 可以从多视角图像生成照片级背景，但仿真器还需要 mesh、碰撞体、质量、摩擦、关节、可变形物体和机器人坐标标定。因此这些方法的差别不在于是否用了 3DGS，而在于它们如何把高斯外观和物理世界接起来。



% METHOD: 2503.00370
\paragraph{Scalable Real2Sim（机器人交互式资产生成）（2025.03）}
Scalable Real2Sim~\cite{2503_00370} 是本章 Real-to-Sim 路线的代表性锚点。原始 3DGS 或 NeRF 可以把物体外观重建得很细，但机器人仿真还需要可碰撞几何、质量、质心和惯性；这些量如果仍靠人工测量，数字孪生就很难规模化。Scalable Real2Sim 的改变是把机器人本身变成自动物理扫描器：机械臂在标准 pick-and-place 设置中反复抓取、放置和翻动物体，外部 RGB-D 相机记录外观，关节扭矩传感器记录物理响应，最后输出可直接进入仿真的视觉网格、碰撞几何和惯性参数。

具体流程分成几步。视觉侧先用视频分割去掉背景和夹爪像素，并通过物体跟踪估计相机坐标系下的物体位姿；机器人会重新抓取物体以暴露不同表面。随后方法从 3DGS 或 NeRF 这类光度重建中提取物体中心 mesh，但不直接把背景 mask 当作硬裁剪，而是采用 alpha 透明训练：背景像素被替换为与迭代相关的随机颜色，迫使模型在背景区域分配低密度，从而减少密度渗入和遮挡残留。物理侧则让机器人执行预设计的激励轨迹，利用空载和持物两种状态下的关节扭矩差异估计物体质量、质心和转动惯量；几何网格再经过凸分解生成碰撞体。公开实现也体现了这种拆分：数据采集入口负责让机器人连续取放并保存 RGB-D、mask 和扭矩，资产生成入口再调用 BundleSDF、Frosting 或 Neuralangelo 等重建模块，最后输出仿真器可用的几何和物理文件。相对“用 3DGS 做一个好看的物体模型”，它真正增量在于把外观重建、物体分割、位姿跟踪和机器人动力学辨识串成一条无需人工标注的资产生成管线。

\textbf{实现细节：} Scalable Real2Sim（2025.03） 提出全自动 Real2Sim 管线，仅使用机器人关节扭矩传感器和外部相机生成仿真就绪资产。与高保真数字孪生方法依赖外部分割模型和语义处理管线不同，Scalable Real2Sim 改变了资产获取的数据来源和物理参数辨识方式——利用机器人自身执行预设计的激励轨迹（有限傅里叶级数参数化），通过半定规划（SDP）从关节扭矩和运动学状态同时辨识机器人连杆的惯性参数（质量、质心、转动惯量）和关节摩擦系数，然后通过差分法从空载和持物两种状态下的参数差获取物体的完整惯性属性。设计动机在于：传统仿真资产生成需要人工标注几何和物理属性，无法规模化；而机器人自身的关节扭矩传感器可以作为物理参数辨识的"天然标定工具"，无需额外传感器。在几何重建方面，alpha 透明训练方法将背景像素替换为迭代相关的随机颜色，迫使模型在背景区域分配零密度，解决了以物体为中心重建中的密度渗透问题；DINO 特征空间的帧选择策略确保下采样后的训练图像保持最大多样性。最终生成的资产包括纹理网格、凸分解碰撞几何和物理属性，可直接集成至现有拾取-放置系统。
\textbf{实现补充：} 实现上还要核对输入、输出、参数、特征、初始化、训练损失、正则、渲染或推理接口，以及优化中的剪枝、分裂、采样、编码或解码步骤；这些抓手决定该方法到底改了 3DGS 管线的哪一环。
\textbf{锚点细化：} Scalable Real2Sim：提出全自动 Real2Sim 管线，仅使用机器人关节扭矩传感器和外部相机生成仿真就绪资产。与高保真数字孪生方法依赖外部分割模型和语义处理管线不同，Scalable Real2Sim 改变了资产获取的数据来源和物理参数辨识方式——利用机器人自身执行预设计的激励轨迹（有限傅里叶级数参数化），通过半定规划（SDP）从关节扭矩和运动学状态同时辨识机器人连杆的惯性参数（质量、质心、转动惯量）和关节摩擦系数，然后通过差分法从空载和持物两种状态下的参数差获取物体的完整惯性属性。设计动机在于：传统仿真资产生成需要人工标注几何和物理属性，无法规模化；而机器人自身的关节扭矩传感器可以作为物理参数辨识的``天然标定工具''，无需额外传感器。在几何重建方面，alpha 透明训练方法将背景像素替换为迭代相关的随机颜色，迫使模型在背景区域分配零密度，解决了以物体为中心重建中的密度渗透问题；DINO 特征空间的帧选择策略确保下采样后的训练图像保持最大多样性。最终生成的资产包括纹理网格、凸分解碰撞几何和物理属性，可直接集成至现有拾取-放置系统。
\textbf{方法信息：} Amazon 参与了这篇论文；代码已公开：\url{https://github.com/nepfaff/scalable-real2sim}，仓库约 137 stars。



% METHOD: 2504.13175
\paragraph{RoboSplat（单示范操控的数据增强）（2025.04）}
Novel Demonstration Generation with Gaussian Splatting Enables Robust One-Shot Manipulation~\cite{2504_13175} 把 3DGS 当作数据增强器。单样本机器人示范的主要问题是策略只见过一个物体、一个视角和一种光照；RoboSplat 的改变是通过编辑 3DGS 重建的场景生成多样化示范，而不重新采集真实数据。

方法先重建包含机器人和环境的 3DGS，并通过 COLMAP 和深度先验对齐坐标。增强包括五类：用 3D 高斯替换改变物体类型、场景外观和机器人形态；用等变变换改变物体姿态；用视觉属性编辑模拟光照；用新视角合成改变相机；用 3D 内容生成增加物体多样性。为了让增强帧仍对齐真实动作，方法加入可微渲染的精细帧对齐。相对 GenSplat 的相机位姿增强，RoboSplat 更全面地编辑示范内容本身。
\textbf{实现细节：} 基于 DMP 的示范合成（2026.05） 将专家轨迹用动态运动基元（DMP）建模并重定向至新目标、物体配置和视角，产生阶段一致、形状保持的运动。分析式障碍感知 DMP 公式直接在 3DGS 诱导的连续密度场上操作实现碰撞避免。与 SplatSim 通过仿真器正向运动学驱动高斯渲染不同，RoboSplat 改变了帧对齐和示范生成的精度——利用 3DGS 可微渲染替代传统 ICP 进行精细帧对齐，通过最小化渲染分割掩码与 URDF 掩码的像素差异优化变换矩阵，解决高斯基元尺度属性导致的点云坐标与外观不匹配问题。设计动机在于：传统 ICP 仅基于点坐标对齐，忽略了 3DGS 特有的尺度属性导致的系统性偏移。RoboSplat 覆盖六种泛化维度（物体位姿、物体类型、相机视角、机器人本体、场景外观、光照条件）的系统化增强策略中，物体类型增强通过 GPT-4 生成可抓取物体名称并使用 3D 内容生成模型创建高斯资产，场景外观增强利用 COCO 数据集和重建 3DGS 场景替换背景，光照增强通过随机缩放、偏移和高斯噪声调制每个高斯的漫反射颜色。
\textbf{方法信息：} 引用很多，Semantic Scholar 约 69 次。
\textbf{设计动机与边界：} 相对原始 3DGS 以独立高斯和图像重建误差为主的设定，Novel Demonstration Generation with Gaussian Splatting Enables Robust One-Shot Manipulatio 把问题落在以下可实现的改动上：RoboSplat首先通过3DGS重建完整场景（含机器人），利用COLMAP获取相机姿态和深度先验。关键创新在于将3DGS作为一种可编辑的3D表示，实现了五种数据增强技术：(1)3D高斯替换用于改变物体类型、场景外观和机器人形态；(2)等变变换用于改变物体姿态；(3)视觉属性编辑用于模拟不同光照条件；(4)新视角合成用于生成不同相机视角；(5)3D内容生成用于增加物体多样性。还设计了基于可微渲染的精细帧对齐方法，解决了3DGS坐标与真实世界坐标系的对齐问题。其设计目的不是单独增加一个模块，而是用上述表示、监督或渲染约束消除原流程在目标任务中的具体失配。提出RoboSplat方法，通过直接编辑3DGS重建的场景生成多样化、视觉真实的机器人操作演示，涵盖物体类型、场景外观、机器人形态、物体姿态、光照条件和相机视角六种泛化类型，在单样本设置下将策略成功率从57.2%提升至87.8%。适用时还应保留论文给出的边界：依赖3DGS重建质量；对透明或反光物体的重建效果有限；需要精确的坐标对齐。
% METHOD: 2505.14938
\paragraph{Scan, Materialize, Simulate（扫描-材质化-仿真）（2025.05）}
Scan, Materialize, Simulate~\cite{2505_14938} 把 3DGS、视觉基础模型、VLM 材料推理和物理仿真组合成物体中心规划框架。它不重新学习物理动力学，而是把视觉重建转成仿真可执行对象，再在物理引擎中搜索动作。

Scan 阶段用多视角 RGB-D 和 3DGS 重建详细几何，同时用视觉基础模型做物体检测与分割，把高斯或体积基元关联到对象。Materialize 阶段把这些基元转换成仿真网格，并让 VLM 根据外观和语义推理材料属性，如密度、摩擦和弹性。Simulate 阶段把对象实例化到物理引擎，优化候选动作再执行。相对 Scalable Real2Sim 的传感器式物性估计，SMS 更强调用 VLM 语义常识补材料属性。
\textbf{实现细节：} 数据增强与合成数据的范式分化。沿 GenSplat 视角增强思路，研究界在"高斯基元操控 $\to$ 轨迹重定向 $\to$ 端到端引擎"三个层次分化。轨迹重定向方向，R2R2R 仅需智能手机物体扫描与单个人类示范视频即可渲染数千个高视觉保真、机器人无关的示范数据，先用 3DGS 重建物体资产再通过 GARField 分割转换为纹理网格；基于 DMP 的示范合成 将专家轨迹用动态运动基元建模并重定向至新目标、物体配置和视角，直接在 3DGS 诱导的连续密度场上操作、实现碰撞避免；RoboPearls 基于增量语义蒸馏和 3D 正则化 NNFM 损失支持多种场景操控操作并配合 LLM 自动化仿真制作。端到端引擎方向，基于 3DGS 的合成数据生成 在 Unreal Engine 中利用照片级高斯泼溅快速生成标注合成数据；SMS 集成 3DGS 场景重建、视觉基础模型语义分割、VLM 材质属性推理和物理仿真，实现可泛化的物理推理和以物体为中心的规划。
\textbf{方法信息：} NVIDIA 参与了这篇论文。
\textbf{设计动机与边界：} 相对原始 3DGS 以独立高斯和图像重建误差为主的设定，Scan, Materialize, Simulate 把问题落在以下可实现的改动上：SMS框架包含三个阶段：（1）Scan：从多视角RGBD图像使用3DGS进行详细几何重建，同时利用视觉基础模型进行物体检测和分割，将体积基元与对应物体关联。（2）Materialize：将体积基元转换为物理仿真可用的网格，利用视觉语言模型（VLM）推理材料属性（密度、摩擦系数等）。从外观和语义信息估计物理特性。（3）Simulate：将重建的网格实例化到物理仿真器中，优化候选机器人动作后再在目标环境中执行。框架的关键创新在于将可微渲染、语义理解和物理仿真桥接，实现跨场景的通用物理推理，无需重新学习基础物理动力学。其设计目的不是单独增加一个模块，而是用上述表示、监督或渲染约束消除原流程在目标任务中的具体失配。提出了SMS统一框架，将3DGS场景重建、视觉基础模型语义分割、VLM材料属性推理和物理仿真整合，实现了无需重新学习物理动力学的通用物理推理和以物体为中心的规划。在台球操控和四旋翼着陆任务上验证了框架的有效性。适用时还应保留论文给出的边界：3DGS重建质量影响后续物理推理精度；VLM材料推理可能不够精确；物理仿真的保真度受网格质量限制。
% METHOD: 2506.04120
\paragraph{Splatting Physical Scenes（SplatMesh 端到端 Real-to-Sim）（2025.06）}
Splatting Physical Scenes~\cite{2506_04120} 用 SplatMesh 混合表示把外观和物理几何紧耦合。3DGS 高保真但不适合碰撞；mesh 可仿真但外观不够细。SplatMesh 把高斯约束在三角面上，用重心坐标确定高斯位置，mesh 变形时高斯随动，从而让同一个资产既能可微渲染又能可微物理。

训练中，MuJoCo MJX 提供可微物理，3DGS 提供可微渲染，系统从不完美机器人轨迹和 RGB 图像中同时优化场景几何、外观、机器人位姿和物理参数。几何通过固定拓扑 mesh 的顶点位移更新，高斯用 surfel 约束贴合法线方向。它还支持从球体开始的渐进式 mesh 初始化和自动分割。相对 Scalable Real2Sim 的“先重建资产再估计物理属性”，Splatting Physical Scenes 更强调端到端联合优化。
\textbf{方法信息：} Google / DeepMind 参与了这篇论文。
\textbf{设计动机与边界：} 相对原始 3DGS 以独立高斯和图像重建误差为主的设定，Splatting Physical Scenes 把问题落在以下可实现的改动上：核心方法是SplatMesh，一种将三角形网格（几何）与3D高斯（外观）紧耦合的混合表示。高斯被约束在网格面上，通过重心坐标确定位置，网格变形时高斯随动。几何通过顶点位移优化，保持固定拓扑连接。端到端优化利用MuJoCo MJX的可微物理模拟和3DGS的可微渲染，从RGB图像最小化光度误差，同时优化场景几何、外观、机器人位姿和物理参数。引入Surfel约束使高斯沿面法线对齐。支持渐进式mesh初始化（从球体开始）和自动分割。其设计目的不是单独增加一个模块，而是用上述表示、监督或渲染约束消除原流程在目标任务中的具体失配。提出SplatMesh混合表示，将3DGS外观与显式网格几何结合，通过MuJoCo MJX实现可微物理与可微渲染的端到端联合优化，可从不完美的机器人轨迹数据直接生成物理仿真可用的场景资产，同时完成物体网格重建、新视角合成和无标注机器人位姿标定。适用时还应保留论文给出的边界：依赖初始粗略网格（球体）；对大型复杂场景的扩展性未验证；物体分割需要额外先验；物理参数优化受限于MuJoCo模拟能力。
% METHOD: 2506.22756
\paragraph{RoboPearls（可编辑视频仿真）（2025.06）}
RoboPearls~\cite{2506_22756} 从演示视频构建可编辑的机器人操作仿真。它的锚点是动态 3DGS 与视频仿真：真实演示数据采集昂贵，若能把一段视频变成可编辑仿真，就能替换物体、增删物体和生成新的训练情境。

方法扩展 3DGS 以支持时间传播，捕获场景时空动态，并从 SAM 增量蒸馏语义特征，使场景能被对象级理解。仿真操作符由增量语义蒸馏和 3D 正则化最近邻特征匹配损失驱动，支持物体替换、添加和删除。LLM 用于解释用户命令并自动组合操作，VLM 用于分析机器人学习失败并形成闭环改进。相对 RoboSplat 的离线示范增强，RoboPearls 更像可编辑视频仿真生产线。
\textbf{实现细节：} RoboPearls（2025.06） 基于增量语义蒸馏和 3D 正则化 NNFM 损失支持多种场景操控操作，结合 LLM 自动化仿真制作和 VLM 分析关闭仿真回路。与 R2R2R 将 3DGS 转换为网格渲染不同，RoboPearls 改变了 3DGS 场景编辑的语义粒度——通过增量语义蒸馏（ISD）机制在运行时按需细化分割粒度，当用户请求的物体粒度（如炉灶上的小按钮）超出 SAM 初始掩码的分辨率时，系统先用 G-DINO 验证当前掩码是否匹配目标，不匹配时用边界框作为 SAM 的精细提示生成新掩码，仅微调相关高斯的身份编码而非重新训练整个场景。设计动机在于：SAM 的多层级掩码无法覆盖所有物体粒度，而重新训练整个场景的语义嵌入成本过高。RoboPearls 还引入 3D 正则化 NNFM 损失限制风格迁移仅作用于目标物体高斯并保持背景不变，多 LLM 智能体协作自动化整个仿真制作流程。
\textbf{实现补充：} 实现上还要核对输入、输出、参数、特征、初始化、训练损失、正则、渲染或推理接口，以及优化中的剪枝、分裂、采样、编码或解码步骤；这些抓手决定该方法到底改了 3DGS 管线的哪一环。
\textbf{方法信息：} ByteDance 参与了这篇论文。



% METHOD: 2510.10637
\paragraph{RoboSimGS（高保真仿真数据生成）（2025.10）}
High-Fidelity Simulated Data Generation for Real-World Zero-Shot Robotic Manipulation Learning with Gaussian Splatting~\cite{2510_10637} 关注零样本 sim-to-real 操控。它把场景分成静态背景、交互物体和机器人三类资产：背景用 3DGS 最大化视觉真实，交互物体用 mesh 保证物理，机器人用标准 URDF。

RoboSimGS 的第一阶段从多视角真实图像构建解耦场景；第二阶段将 3DGS 背景和 mesh 物体导入物理模拟器。它使用多模态大语言模型从视觉数据推断物体物理属性和运动学结构，例如铰链、滑轨、密度或刚度，再通过物体、相机、光照和轨迹随机化生成多样仿真数据。策略只在仿真数据上训练后部署到真实世界。相对 GSWorld 的闭环资产格式，RoboSimGS 更突出 MLLM 自动生成物理合理铰接资产。
\textbf{实现细节：} RoboSimGS（2025.10） 将多视角真实世界图像转换为可扩展、高保真、物理交互的仿真环境。与 SplatSim 依赖手动分割和 ICP 对齐不同，RoboSimGS 改变了资产生成的自动化程度——采用解耦重建策略（3DGS 用于静态背景的高保真渲染、显式网格用于可交互物体的物理仿真），并通过 MLLM（GPT-4o）分析视觉数据推断物体的运动学结构（关节类型、运动轴、运动极限）和物理属性（密度、杨氏模量、泊松比），自动生成可交互的关节物体 URDF 资产。设计动机在于：传统仿真环境构建需要大量人工建模和参数调优，而 MLLM 的常识推理能力可以从多视角图像推断物体的物理和运动学属性，实现从真实照片到物理仿真环境的全自动转换。系统还通过全场景增强（物体位姿/尺度随机化、相机视角扰动、高斯颜色/光照抖动、轨迹中间点随机化）进一步缩小域间隙，实现零样本仿真到真实迁移。
\textbf{方法信息：} Alibaba 参与了这篇论文。
\textbf{设计动机与边界：} 相对原始 3DGS 以独立高斯和图像重建误差为主的设定，High-Fidelity Simulated Data Generation for Real-World Zero-Shot Robotic Manipulation Lear 把问题落在以下可实现的改动上：RoboSimGS是一个两阶段管线。第一阶段为场景重建：采用解耦重建方法，3DGS用于静态背景以最大化视觉保真，显式网格用于交互式物体以确保物理合理性，机器人本身使用标准URDF文件。每个任务场景需约10分钟的手动扫描获取多视角数据。第二阶段为Sim2Real环境对齐：将3DGS模型作为视觉背景、网格作为物理资产导入物理模拟器。首创使用MLLM分析视觉数据，推断物理属性（密度、刚度）和复杂运动学结构（铰链、滑轨等）。还引入全局场景增强（随机化物体、相机、光照和轨迹）来创建多样化仿真数据。策略完全在增强仿真数据上训练后可直接零样本迁移到真实世界。其设计目的不是单独增加一个模块，而是用上述表示、监督或渲染约束消除原流程在目标任务中的具体失配。提出了RoboSimGS框架，将多视角真实世界图像转换为可扩展、高保真且具有物理交互能力的仿真环境。首创使用多模态大语言模型（MLLM）自动创建物理合理的铰接资产，推断物体的物理属性和运动学结构。采用3DGS+Mesh混合表示实现视觉保真与物理模拟的统一。仅在仿真数据上训练的策略即可在真实世界实现零样本迁移。适用时还应保留论文给出的边界：场景重建需要约10分钟的手动扫描；MLLM推断物理属性的准确性受限于模型的物理理解能力；主要验证了桌面级操作任务，大规模场景的适用性待验证。
% METHOD: 2510.15352
\paragraph{GaussGym（3DGS 接入向量化物理仿真器）（2025.10）}
GaussGym~\cite{2510_15352} 的目标是把 3DGS 作为即插即用渲染器嵌入 IsaacGym 等向量化物理仿真器。多数 3DGS real-to-sim 方法重视场景重建，但策略训练还要求并行环境、高吞吐和物理控制接口；GaussGym 的增量在工程集成。

数据管线可从 ARKitScenes、GrandTour、手机扫描或生成式视频模型输出开始，用 VGGT 提取相机参数和稠密表示，用 NKSR 生成 mesh，高斯则从 VGGT 点云初始化以加速收敛。渲染管线在大量并行环境中 rasterize splats，并将渲染频率和控制频率解耦，按相机真实帧率渲染而不是每个控制步都渲染。它用辅助重建损失利用真实 mesh 数据加快学习，面向人形和四足运动策略训练。相对 GSWorld，它更偏开源仿真后端与高吞吐训练。
\textbf{方法信息：} 引用很多，Semantic Scholar 约 25 次。
\textbf{设计动机与边界：} 相对原始 3DGS 以独立高斯和图像重建误差为主的设定，GaussGym 把问题落在以下可实现的改动上：GaussGym的核心是将3DGS作为即插即用渲染器嵌入向量化物理仿真器。数据管线：支持posed数据集（ARKitScenes、GrandTour）、智能手机扫描、甚至生成式视频模型（Veo）的输出。使用VGGT提取相机参数和密集场景表示，NKSR生成高质量网格，高斯泼溅直接从VGGT点云初始化以加速收敛。渲染管线：在4096个环境中并行光栅化泼溅，使用多线程PyTorch内核确保高效GPU利用。关键优化：将渲染与控制频率解耦，按相机真实帧率渲染而非控制频率；使用辅助重建损失利用真实网格数据加速学习。支持人形和四足机器人的运动导航策略训练。其设计目的不是单独增加一个模块，而是用上述表示、监督或渲染约束消除原流程在目标任务中的具体失配。提出开源real-to-sim框架GaussGym，将3DGS作为即插即用渲染器集成到向量化物理仿真器（如IsaacGym）中，在消费级GPU上实现超过100,000步/秒的仿真速度和照片级渲染质量，支持iPhone扫描、大规模数据集和生成式视频模型输出的多样化环境输入。适用时还应保留论文给出的边界：RGB策略学习仍具挑战性，需要辅助损失加速；GS渲染不提供精确碰撞检测；sim-to-real迁移仍有差距。
% METHOD: 2511.04665
\paragraph{Soft-Body Real-to-Sim Policy Evaluation（可变形物体策略评估）（2025.11）}
Real-to-Sim Robot Policy Evaluation with Gaussian Splatting Simulation of Soft-Body Interactions~\cite{2511_04665} 关注软体物体操控。普通 3DGS 数字孪生能渲染真实外观，但软体任务的策略评估还需要物体在按压、抓取和拉扯时的动力学。该方法把 PhysTwin 的弹簧-质点系统与 3DGS 外观模型结合起来。

动力学部分从人手交互视频中优化弹簧阈值和刚度等物理参数，把可变形物体表示为密集弹簧-质点系统。外观部分用 3DGS 从单次扫描重建场景，提供任意视角光真实渲染，并加入自动位姿、颜色对齐和物体变形处理，使渲染能跟随动力学状态变化。二者被封装进 Gym 风格接口，用于在真实部署前评估机器人策略。相对 Scalable Real2Sim 的刚体资产生成，它的重点是软体动力学和策略评估。
\textbf{实现细节：} SoMA（2026.02） 提出面向软体操控的 3DGS 仿真器，在统一潜在神经空间耦合可变形动力学、环境力和机器人关节动作，实现端到端 Real-to-Sim 仿真。与 GaussTwin 通过 PBD 物理引擎驱动高斯基元不同，SoMA 改变了软体动力学建模的方式——采用层次图结构建模机器人-物体-环境交互，将环境效果建模为显式全局力（重力和支持力），机器人交互通过交互图中的隐式力计算，两者直接作用于高斯基元层级。设计动机在于：现有的神经高斯动力学模型主要基于状态回归，将物体运动与接触效果纠缠在一起，在多样化交互模式下难以泛化。SoMA 通过力驱动的动力学解耦（$\dot{v}_i, \dot{\omega}_i = f_\theta(g_i^{t-1}, g_i^{t-2}, f_i)$，其中 $f_i = f_i^{}\mathrm{env} + f_i^{}\mathrm{rob}$），使模型学习力到状态的映射而非直接的状态转移，从而在未见过的交互模式中保持泛化能力。
\textbf{实现补充：} 实现上还要核对输入数据、输出指标、基线配置、评测协议、系统接口和适用边界；若调用 3DGS 渲染或重建结果，也应说明相机参数、训练设置、优化流程和指标计算如何对齐。
\textbf{设计动机与边界：} 相对原始 3DGS 以独立高斯和图像重建误差为主的设定，Real-to-Sim Robot Policy Evaluation with Gaussian Splatting Simulation of Soft-Body Intera 把问题落在以下可实现的改动上：框架包含两个核心组件：(1)动力学模型：采用PhysTwin框架，将可变形物体重建为密集弹簧-质点系统，从人手交互视频优化物理参数（弹簧阈值d和刚度系数Y）；(2)外观模型：利用3DGS从单次扫描重建光真实感场景，支持任意视角渲染，并增强自动位姿和颜色对齐以及物体变形处理能力。两者集成于Gym风格接口的统一仿真器中。其设计目的不是单独增加一个模块，而是用上述表示、监督或渲染约束消除原流程在目标任务中的具体失配。提出Real-to-Sim策略评估框架，结合PhysTwin的弹簧-质点系统物理引擎和3DGS光真实感渲染，从真实世界视频构建可变形物体的数字孪生。实现了机器人、物体和环境的自动位姿与颜色对齐以及物体变形处理，弥合仿真与真实之间的外观和动力学差距。适用时还应保留论文给出的边界：弹簧-质点模型对高度非线性变形材料建模能力有限；3DGS渲染的新视角合成在训练数据不足区域可能出现伪影；需要手动进行初始场景设置。
\textbf{方法信息：} Google / DeepMind 参与了这篇论文；引用很多，Semantic Scholar 约 23 次。



% METHOD: 2601.03200
\paragraph{3DGS Robotic Digital Twin（规划就绪数字孪生）（2026.01）}
A High-Fidelity Digital Twin for Robotic Manipulation Based on 3D Gaussian Splatting~\cite{2601_03200} 面向可直接接入 Unity、ROS2 和 MoveIt 的操作规划。它不是只输出可视化 3DGS，而是把高斯转成有语义标签和碰撞几何的数字孪生。

几何流采用 InstantSplat 和 MASt3R 先验，从稀疏 RGB 输入直接估计点云和相机位姿，再自监督优化 3DGS，绕过传统 SfM。语义流提出可见性感知语义融合：把 2D 分割掩码当作噪声空间假设，通过深度聚类隔离对象，并用可见性加权投影在多视图聚合标签，减轻语义 bleeding。最后通过统计滤波、Alpha Shape 表面重建、网格简化和碰撞模型提取，把高斯输出转成规划就绪几何。相对 GaussTwin，它更强调分钟级构建和机械臂规划接口。
\textbf{方法信息：} 代码：\url{https://github.com/535a59/3dgs-digital-twin}。
\textbf{设计动机与边界：} 相对原始 3DGS 以独立高斯和图像重建误差为主的设定，A High-Fidelity Digital Twin for Robotic Manipulation Based on 3D Gaussian Splatting 把问题落在以下可实现的改动上：框架包含两条并行处理流：(1)几何重建流：采用InstantSplat方法利用预训练几何先验(MASt3R)直接估计初始点云和相机位姿，通过自监督优化联合优化3DGS参数和相机位姿，绕过传统SfM步骤实现快速收敛；(2)语义分割流：提出可见性感知语义融合框架，将2D分割掩码视为噪声空间假设，通过深度聚类实现空间隔离确保语义一致性，利用可见性加权投影从多视图聚合语义标签，解决'bleeding effect'问题。随后通过滤波式几何转换方法：基于统计滤波的点云去噪、Alpha Shape表面重建、网格简化与碰撞模型提取，将3DGS输出转换为规划就绪的碰撞几何。最终与Unity-ROS2-MoveIt集成，支持Franka Emika Panda机器人执行pick-and-place任务。其设计目的不是单独增加一个模块，而是用上述表示、监督或渲染约束消除原流程在目标任务中的具体失配。提出基于3DGS的快速高保真数字孪生框架，通过可见性感知语义融合实现精确3D标注，引入高效滤波式几何转换方法生成碰撞就绪模型，与Unity-ROS2-MoveIt物理引擎无缝集成，在分钟级别从稀疏RGB输入构建支持机器人操控的高质量数字孪生。适用时还应保留论文给出的边界：3DGS输出的几何模糊性（漂浮物、鬼影伪影、表面模糊）需要额外的滤波处理；语义融合依赖2D分割模型质量。
% METHOD: 2603.05108
\paragraph{GaussTwin（预测-校正数字孪生）（2026.03）}
GaussTwin~\cite{2603_05108} 把机器人数字孪生做成预测-校正闭环。单纯仿真会随着真实世界扰动逐渐偏离，单纯视觉跟踪又缺少物理预测；GaussTwin 的改变是把高斯绑定到物理基元上，用物理前推状态，再用视觉误差校正。

物理部分使用基于位置的动力学处理刚体碰撞和接触，用离散 Cosserat rod 公式处理绳索等弹性杆件。初始化时从多相机 RGB-D 构建物体高斯，每个粒子或刚体绑定一组 Gaussian 用于渲染。运行时以约 25Hz 进行预测-校正：预测步由物理仿真前推，校正步用光度误差和分割掩码驱动 SE(3) 更新，把仿真状态拉回真实观测。相对 Splatting Physical Scenes 的端到端建模，它更强调在线同步和可变形体闭环。
\textbf{实现细节：} GaussTwin（2026.03） 将基于位置的动力学（PBD）与离散 Cosserat 杆公式结合实现物理接地仿真。与高保真数字孪生方法的静态重建-网格提取管线不同，GaussTwin 改变了 3DGS 的动态更新环节——通过将高斯锚定至物理基元（球体粒子）并强制一致 SE(3) 更新，实现 25\,Hz 实时的预测-校正闭环。设计动机在于：刚性物体和可变形体（如绳索）的操控需要同时具备物理一致性（碰撞、接触力）和视觉真实性，而传统方法中物理仿真和视觉渲染使用独立表示，无法实现两者的闭环耦合。对象-高斯绑定。场景初始化阶段，GaussTwin 从多视角 RGB-D 图像中用 SAM2 提取实例掩码，在物体包围盒内均匀填充半径 5\,mm 的球体粒子，通过深度投影剔除掩码外、地面以下或深度不一致的粒子，剩余粒子近似物体几何。由粒子密度积分计算物体的质量、质心、位姿和惯性矩阵。3D 高斯从分割后的物体点云初始化，经 1000 步 Adam 优化后锚定至最近的球体粒子，建立高斯外观与物理基元的固定绑定关系。
\textbf{方法信息：} NVIDIA 参与了这篇论文。
\textbf{设计动机与边界：} 相对原始 3DGS 以独立高斯和图像重建误差为主的设定，GaussTwin 把问题落在以下可实现的改动上：GaussTwin将物理仿真和视觉表示统一：(1)物理仿真：使用基于位置的动力学(PBD)处理刚体碰撞和接触约束，使用离散Cosserat杆公式建模弹性杆件的剪切-拉伸和弯曲-扭转约束；(2)物体-Gaussian绑定：从多相机RGB-D图像初始化构建，每个物理基元（粒子或刚体）绑定一组Gaussian用于渲染；(3)预测-校正循环：25Hz运行，预测步使用物理仿真前推状态，校正步使用光度误差和分割掩码驱动的SE(3)更新同步仿真与现实。支持刚体和可变形体（绳索等）的统一处理。其设计目的不是单独增加一个模块，而是用上述表示、监督或渲染约束消除原流程在目标任务中的具体失配。提出GaussTwin实时数字孪生，结合基于位置的动力学(PBD)与离散Cosserat杆公式实现物理仿真，3DGS进行高效渲染和视觉校正。通过将Gaussian锚定到物理基元并使用光度误差和分割掩码驱动的SE(3)更新实现稳定的预测-校正。适用时还应保留论文给出的边界：目前针对特定类型的可变形体（杆件）。物理仿真的参数需要针对不同材质进行调整。
% METHOD: 2603.29192
\paragraph{GenSplat（相机位姿增强的策略泛化）（2026.03）}
Efficient Camera Pose Augmentation for View Generalization in Robotic Policy Learning~\cite{2603_29192} 解决策略对相机安装位置敏感的问题。真实机器人数据通常来自固定相机配置，一旦部署时相机外参变化，端到端策略会退化。GenSplat 的改变是从稀疏未标定输入前馈重建 3DGS，再渲染大量合成视角作为训练增强。

框架用 DINOv2 patch token 和交替的帧内/全局自注意力聚合多视图上下文，保持输入视角置换等变。相机位姿解码器预测局部到全局变换，高斯预测器输出位置、协方差、不透明度和颜色。为了避免纯光度训练导致几何坍塌，它从 VGGT 等基础模型蒸馏 3D 先验。策略训练时，从重建高斯中采样多样化相机视角，扩展观察流形。相对 RoboSplat 的内容增强，它专攻相机位姿泛化。
\textbf{实现细节：} GenSplat（2026.03） 提出前馈 3DGS 框架，通过新视角渲染增强视觉运动策略的视角泛化能力。与高保真数字孪生方法将 3DGS 用于场景重建不同，GenSplat 改变了 3DGS 的用途——不重建特定场景，而是为策略训练批量生成视角增强数据。设计动机在于：主流 2D 视觉运动策略对训练时未见过的相机视角严重退化，而大规模机器人数据集通常缺乏可靠的相机位姿标注，无法直接使用传统 3DGS 进行新视角合成。置换等变架构。GenSplat 采用前馈置换等变 Transformer 架构，从稀疏未标定图像直接预测高斯参数和相机位姿，无需逐场景优化。每张图像先经 DINOv2 提取 patch token，通过交替的帧内注意力和全局自注意力层聚合多视角上下文，确保输出与输入视角顺序无关。token 随后分别路由至相机位姿解码器（预测局部到全局的 SE(3) 变换）和高斯参数解码器（预测位置、颜色、不透明度和缩放）。这一架构的关键差异在于：与传统 3DGS 依赖 COLMAP/SfM 逐场景重建不同，GenSplat 在单次前向传播中完成从任意数量未标定视图到完整 3DGS 表示的映射，推理速度比逐场景优化快数个量级。3D 先验蒸馏策略。纯光度监督在稀疏视图、无位姿条件下容易导致几何坍塌和渲染伪影。GenSplat 通过从预训练几何基础模型 PI3D 提取密集点图监督和仿射不变位姿约束，正则化 3DGS 优化。直觉上，PI3D 为场景提供了一个粗略但合理的几何"骨架"，防止优化过程陷入退化的局部最优（如所有高斯退化为平面或零厚度壳体）。视角增强管线与实验。对于每条专家示范，GenSplat 重建 3DGS 场景后，在相机绕世界轴旋转 $[-30^{}\circ, 30^{}\circ]$ 并沿 $+Z$ 轴平移 $[0.5, 6]$\,cm 的范围内采样新视角渲染合成轨迹，严格保持时间对齐和动作标签不变。该策略对策略架构（Diffusion Policy 和 ACT）均有效，实现了架构无关的视角泛化增强。
\textbf{技术实现补充：} GenSplat（2026.03） 提出前馈 3DGS 框架，通过新视角渲染增强视觉运动策略的视角泛化能力。与高保真数字孪生方法将 3DGS 用于场景重建不同，GenSplat 改变了 3DGS 的用途——不重建特定场景，而是为策略训练批量生成视角增强数据。设计动机在于：主流 2D 视觉运动策略对训练时未见过的相机视角严重退化，而大规模机器人数据集通常缺乏可靠的相机位姿标注，无法直接使用传统 3DGS 进行新视角合成。置换等变架构。GenSplat 采用前馈置换等变 Transformer 架构，从稀疏未标定图像直接预测高斯参数和相机位姿，无需逐场景优化。每张图像先经 DINOv2 提取 patch token，通过交替的帧内注意力和全局自注意力层聚合多视角上下文，确保输出与输入视角顺序无关。token 随后分别路由至相机位姿解码器（预测局部到全局的 SE(3) 变换）和高斯参数解码器（预测位置、颜色、不透明度和缩放）。这一架构的关键差异在于：与传统 3DGS 依赖 COLMAP/SfM 逐场景重建不同，GenSplat 在单次前向传播中完成从任意数量未标定视图到完整 3DGS 表示的映射，推理速度比逐场景优化快数个量级。3D 先验蒸馏策略。纯光度监督在稀疏视图、无位姿条件下容易导致几何坍塌和渲染伪影。GenSplat 通过从预训练几何基础模型 PI3D 提取密集点图监督和仿射不变位姿约束，正则化 3DGS 优化。直觉上，PI3D 为场景提供了一个粗略但合理的几何"骨架"，防止优化过程陷入退化的局部最优（如所有高斯退化为平面或零厚度壳体）。视角增强管线与实验。对于每条专家示范，GenSplat 重建 3DGS 场景后，在相机绕世界轴旋转 $[-30^\circ, 30^\circ]$ 并沿 $+Z$ 轴平移 $[0.5, 6]$\,cm 的范围内采样新视角渲染合成轨迹，严格保持时间对齐和动作标签不变。该策略对策略架构（Diffusion Policy 和 ACT）均有效，实现了架构无关的视角泛化增强。
\textbf{方法信息：} Huawei 参与了这篇论文。
\textbf{设计动机与边界：} 相对原始 3DGS 以独立高斯和图像重建误差为主的设定，Efficient Camera Pose Augmentation for View Generalization in Robotic Policy Learning 把问题落在以下可实现的改动上：GenSplat框架包含：(1)问题建模：每个演示包含多相机视角的观测序列、动作和语言指令。目标是无需额外数据采集即可训练对新视角泛化的策略。构建逐帧无位姿相机增强管线。(2)置换等变架构：使用DINOv2生成patch token，通过交替帧内和全局自注意力层聚合多视图上下文，确保表示严格置换等变。相机位姿解码器预测局部-全局变换，高斯预测器输出高斯参数。(3)3D先验蒸馏：从预训练基础模型（如VGGT）提取几何先验，正则化3DGS优化防止纯光度监督下的几何坍塌。(4)策略训练：从重建的3DGS中采样多样化的合成视角增强训练集，使策略在增强数据上学习3D结构 grounded 的动作映射。其设计目的不是单独增加一个模块，而是用上述表示、监督或渲染约束消除原流程在目标任务中的具体失配。提出GenSplat前馈3DGS框架，从稀疏未标定输入在单次前向传播中重建高保真3D场景，通过渲染多样化合成视角系统增强训练观测流形。设计置换等变架构保证鲁棒性，引入3D先验蒸馏策略正则化3DGS优化防止几何坍塌，使策略在严重空间扰动下仍能稳健执行。适用时还应保留论文给出的边界：前馈3DGS重建质量受输入视角数量限制，3D先验蒸馏增加计算开销，需平衡增强视角的多样性。

\section{导航与路径规划}

导航路线的锚点是语义地图和仿真环境。3DGS 的增量在于，它既能作为可渲染记忆支持新视角观察，也能承载开放词汇语义、可导航区域和动态障碍物。与传统 2D/体素地图相比，3DGS 更连续、更真实，但必须补足不确定性、语义查询和碰撞/通行性。
% METHOD: 2406.19811
\paragraph{EgoGaussian（自我中心视频的动态场景理解）（2024.06）}
EgoGaussian~\cite{2406_19811} 从纯 RGB egocentric video 中同时重建静态背景和动态物体。机器人或可穿戴相机常记录人与物体交互，传统 3DGS 假设静态场景会把动态物体拖成伪影。EgoGaussian 利用 3DGS 显式点状表示做空间分割，并用 clip 级在线学习处理时间。

预处理用 COLMAP 估计相机姿态，并用手-物体分割得到动态掩码。视频被分成静态片段和动态片段：静态片段初始化背景和物体形状，动态片段分离前景物体高斯和背景高斯分别训练，并追踪物体刚体运动。为了简化外观，方法禁用视角相关颜色。它相对 ReaDy-Go 的区别在于，EgoGaussian 是从 egocentric 数据理解动态交互，而 ReaDy-Go 是用重建场景生成导航训练数据。
\textbf{实现细节：} EgoGaussian（2024.06） 首次从仅 RGB 自我中心输入同时重建 3D 场景并动态跟踪 3D 物体运动。与本章其他动态 3DGS 方法（如 Physically Embodied GS）依赖外部物理引擎不同，EgoGaussian 改变了 3DGS 的场景分解和优化策略——利用手-物体交互的时序结构将视频分割为静态片段（无物体运动的交互准备/后续阶段）和动态片段（物体被移动的阶段），分别重建背景和动态物体。设计动机在于：自我中心视角下的手-物体交互视频存在大量遮挡和场景不一致性（物体被拿起后原位置被遮挡），直接对全序列进行 3DGS 优化会产生大量浮游伪影。EgoGaussian 利用高斯泼溅的离散特性，首先在静态片段中训练完整场景的 3DGS 并通过相邻动态片段的物体掩码识别动态物体的高斯基元，然后将物体高斯作为整体进行刚体位姿估计（6D 旋转表示 + 平移），在动态片段中逐帧优化物体位姿并同步精化物体形状，最终将背景模型和动态物体模型合成为完整的 4D 重建结果。
\textbf{方法信息：} 代码：\url{https://github.com/zdwww/egogaussian}。
\textbf{设计动机与边界：} 相对原始 3DGS 以独立高斯和图像重建误差为主的设定，EgoGaussian 把问题落在以下可实现的改动上：方法利用3DGS的显式点云式表示使其特别适合空间分割。预处理阶段：使用COLMAP估计相机姿态，通过现成方法获取手-物体分割掩码。将视频分为静态片段（仅相机运动）和动态片段（含物体交互）。静态片段用于初始化静态背景和物体形状，动态片段用于追踪物体运动并逐步精化形状。clip级在线学习管线：按时间顺序处理视频，对每个动态片段分离前景物体高斯和背景高斯分别训练。通过自然交互轨迹的时间分段，逐步重建场景的时间演化。对于刚体运动追踪，利用物体的3D高斯表示在不同帧之间的对应关系。禁用视角相关颜色（SH degree设为0）以简化模型。整个流程完全基于RGB输入，不需要深度相机或多相机设置。其设计目的不是单独增加一个模块，而是用上述表示、监督或渲染约束消除原流程在目标任务中的具体失配。首个从纯RGB自我中心视频同时重建3D场景和追踪动态物体运动的方法。利用高斯溅射的离散特性实现空间分割，clip级在线学习管线处理时间维度，在动态物体和背景重建质量上显著优于SOTA。适用时还应保留论文给出的边界：主要针对刚体运动，非刚性变形处理有限。依赖手-物体分割质量。对快速运动和严重遮挡场景有挑战。
% METHOD: 2507.01125
\paragraph{VISTA（任务相关开放词汇主动探索）（2025.07）}
VISTA~\cite{2507_01125} 面向机器人主动探索。它在线构建语义 3DGS 地图，并根据开放词汇搜索指令规划信息轨迹。普通 NBV 只看几何覆盖，可能探索与任务无关的区域；VISTA 的改变是同时估计几何视角多样性和语义任务相关性。

在线建图时，系统从 CLIP 等视觉语言模型提取语义特征并增量嵌入 3DGS，形成任务相关热力图。视角-语义覆盖度量中，几何不确定性由候选视角与历史视角之间的最小角度间隔表示，并考虑遮挡；语义相关性由高斯语义特征与查询文本相似度表示。退缩地平线规划采样候选轨迹，最大化几何覆盖和语义相关性的加权目标。相对 LagMemo 的长期记忆查询，VISTA 强调边探索边找目标。
\textbf{实现细节：} ATLAS Navigator 展示了语言嵌入 3DGS 地图在导航中的威力，但有两个关键局限：其一，它假设环境已经预先建图，无法在机器人首次进入的未知环境中运行；其二，它不量化地图的几何不确定性——即无法判断哪些区域尚未被充分观测、哪些区域还需要更多视角来完善重建。当目标物体不在初始视野范围内时，仅靠语义搜索而无几何引导可能导致探索效率急剧下降。VISTA（2025.07） 正是在这一背景下提出的增量方法：它继承了 ATLAS Navigator 的语言嵌入建图思路（将 CLIP 语义特征蒸馏到 3DGS 中），但将整个范式从"在已有地图上搜索"推进到"边建图边搜索"——在未知环境中同步完成语义建图和目标导向的主动探索。在线语义建图。VISTA 基于 SplatBridge（NerfBridge 的 3DGS 扩展）实现实时 3DGS 训练：机器人采集的 RGB-D 图像和位姿持续汇入流式数据集，3DGS 参数同步优化。
\textbf{方法信息：} 代码：\url{https://github.com/fierrojuanpablorenovado-commits/vista3d}。
\textbf{设计动机与边界：} 相对原始 3DGS 以独立高斯和图像重建误差为主的设定，VISTA 把问题落在以下可实现的改动上：VISTA包含三个关键组件：(1) 在线语义3DGS建图：实时构建场景的高保真3DGS表示，同时从视觉语言模型（CLIP）提取语义特征增量嵌入3DGS地图，生成任务相关性3D热力图；(2) 高效视角-语义覆盖度量：同时量化几何视角多样性和任务相关性。在每个体素中，几何不确定性是测试视角与历史所有视角之间的最小角度间隔（考虑遮挡），支持递归更新。语义相关性通过CLIP特征与查询的相似度计算；(3) 退缩地平线轨迹规划：采样多条轨迹，选择视角最大化几何不确定性和语义相关性的加权组合，引导机器人向指定查询物体移动。方法无需预建地图，支持开放词汇搜索指令。其设计目的不是单独增加一个模块，而是用上述表示、监督或渲染约束消除原流程在目标任务中的具体失配。提出VISTA主动探索方法，实时构建语义3DGS地图并基于开放词汇搜索指令规划信息轨迹。核心创新是高效的视角-语义覆盖度量，同时量化几何视角多样性和任务相关性，在四旋翼硬件实验中挑战性地图成功率提升6倍，支持四旋翼和四足机器人。适用时还应保留论文给出的边界：语义特征提取增加计算开销；轨迹规划受限于采样数量
% METHOD: 2509.17430
\paragraph{EmbodiedSplat（手机捕获场景的个性化导航）（2025.09）}
EmbodiedSplat~\cite{2509_17430} 的目标是让普通用户用手机捕获部署环境，再训练个性化 ImageNav 策略。它把 3DGS 重建转成 Habitat-Sim 可用场景，缩小大规模通用数据集和真实家居环境之间的差距。

流程先用 iPhone 和 Polycam 采集 RGB-D 数据，通过 Nerfstudio 获得约千帧对齐图像和位姿，再用 DN-Splatter 重建高质量网格。网格被导入 Habitat-Sim 作为 ImageNav 训练环境，策略先在 HM3D/HSSD 上预训练，再在个性化重建场景上微调。实验显示仿真表现和真实导航有较高相关性。相对 NavGSim 的大规模通用仿真器，EmbodiedSplat 更强调低成本个性化真实场景。
\textbf{实现细节：} 上述方法（GSMem、GLMap 等）均假设策略已在仿真中训练完成，3DGS 的角色局限于在线建图与语义推理——即 3DGS 在导航过程中作为场景表示被构建和消费，但不参与策略本身的训练。EmbodiedSplat（2025.09） 则将 3DGS 的角色从导航阶段的运行时表示前置为训练阶段的环境生成器，提出一个互补视角：能否利用 3DGS 的低成本场景重建能力，直接弥合导航策略的仿真-真实鸿沟？该问题的核心瓶颈在于，传统高质量仿真场景（如 HM3D 的 800 训练场景和 HSSD 的 134 训练场景）依赖 Matterport 扫描仪采集，设备昂贵且耗时数小时，且场景风格偏重公寓布局——当部署环境为大学教室、会议室等分布外场景时，预训练策略的零样本性能急剧下降。EmbodiedSplat 正是针对这一分布外部署难题，提出用消费级手机快速重建部署环境并在其中个性化微调导航策略。Real-to-Sim-to-Real 管线。EmbodiedSplat 的核心思路是用消费级手机快速重建部署环境，在重建场景中微调导航策略后部署到真实机器人。
\textbf{实现补充：} 实现上还要核对输入、输出、参数、特征、初始化、训练损失、正则、渲染或推理接口，以及优化中的剪枝、分裂、采样、编码或解码步骤；这些抓手决定该方法到底改了 3DGS 管线的哪一环。
\textbf{方法信息：} Waymo 参与了这篇论文。
\textbf{设计动机与边界：} 相对原始 3DGS 以独立高斯和图像重建误差为主的设定，EmbodiedSplat 把问题落在以下可实现的改动上：EmbodiedSplat的pipeline：(1) 使用iPhone和Polycam应用捕获部署场景的RGB-D数据（20-30分钟），通过Nerfstudio处理获得1000帧对齐的RGB-D图像和位姿；(2) 使用DN-Splatter（深度-法线正则化的GS方法）重建高质量网格，比Polycam默认网格质量更高；(3) 将重建网格集成到Habitat-Sim中，创建ImageNav训练场景；(4) 在HM3D/HSSD上预训练ImageNav策略，然后在重建场景上微调。实验分析了不同训练策略、预训练数据集和网格重建技术的影响。其设计目的不是单独增加一个模块，而是用上述表示、监督或渲染约束消除原流程在目标任务中的具体失配。提出EmbodiedSplat，首次探索用iPhone捕获的3DGS重建网格训练端到端ImageNav策略并迁移到真实世界。通过GS重建+Habitat-Sim仿真器的pipeline，个性化策略训练，在真实ImageNav任务上相比大规模数据集预训练基线提升20-40%绝对成功率，sim-vs-real相关性达0.87-0.97。适用时还应保留论文给出的边界：iPhone捕获仍需20-30分钟；GS重建质量受捕获质量影响；仅验证了室内导航场景。
% METHOD: 2510.23205
\paragraph{VR-Drive（前馈 3DGS 辅助的视角鲁棒驾驶）（2025.10）}
VR-Drive~\cite{2510_23205} 虽属自动驾驶，但在本章可看作移动机器人视角泛化。它用前馈 3DGS 重建作为辅助任务，让端到端驾驶策略对相机外参变化更鲁棒。普通驾驶网络学到的是固定相机布局下的图像特征；VR-Drive 让网络先预测可从新视角渲染的高斯基元。

原始视角学习阶段，ResNet50 提取多视角特征，像素级预测 3DGS 参数并联合学习深度。新视角学习阶段，随机采样相机外参，用高斯从任意视角渲染特征图，构建视角混合记忆库，促进多视角时序交互；视角一致蒸馏把原始视角知识迁移到合成视角。最终感知和规划端到端训练，缓解合成噪声。相对 GenSplat，它直接把视角增强接入驾驶控制任务。
\textbf{方法信息：} 代码：\url{https://github.com/mickeykang16/VR-Drive}。
\textbf{设计动机与边界：} 相对原始 3DGS 以独立高斯和图像重建误差为主的设定，VR-Drive 把问题落在以下可实现的改动上：VR-Drive框架包含三个组件：(1) 原始视角学习——使用ResNet50提取多视角特征，以像素级前馈方式直接从图像预测3DGS基元（位置、协方差、不透明度、SH颜色），联合学习深度估计，实现无需场景特定优化的在线推理；(2) 新视角学习——随机采样相机外参，利用高斯基元从任意视角渲染多视角特征图，构建视角混合记忆库促进多视角间的时序交互，通过视角一致蒸馏策略将原始视角知识迁移到合成视角；(3) 感知-规划学习——全端到端训练，有效缓解合成引起的噪声。发布新基准数据集评估E2E-AD在新型相机视角下的性能。其设计目的不是单独增加一个模块，而是用上述表示、监督或渲染约束消除原流程在目标任务中的具体失配。提出VR-Drive，首个通过前馈3DGS联合学习3D场景重建作为辅助任务实现视角泛化的端到端驾驶框架，引入视角混合记忆库和视角一致蒸馏策略，有效解决不同车辆配置下相机视角变化导致的鲁棒性问题。适用时还应保留论文给出的边界：前馈3DGS渲染质量可能不如优化式方法；对极端视角变化的泛化能力有限；需要深度和位姿信息。
% METHOD: 2510.24118
\paragraph{LagMemo（语言 3DGS 记忆）（2025.10）}
LagMemo~\cite{2510_24118} 把 3DGS 做成多模态开放词汇导航记忆。一次探索阶段中，机器人用前沿探索收集稀疏 RGB-D 和位姿，重建参数化 3D 高斯，并为每个高斯扩展码本语言特征嵌入。这样目标可以是类别、图像或文本描述，而不仅是固定标签。

它还引入关键帧检索策略：根据 PSNR 负相关权重采样历史帧参与优化，缓解大规模多房间环境中的遗忘和表面孔洞。导航阶段，系统查询语言 3DGS 记忆预测候选目标位置，再用局部感知验证机制动态匹配和验证目标。相对 2D 语义地图，LagMemo 的优势是稀疏视角插值和连续可微优化；相对 VISTA，它更像一次构建后用于多目标导航的长期记忆。
\textbf{实现细节：} 3DGSNav 的虚拟视角重验证证明了 3DGS 渲染对导航的价值，但其记忆服务于单次 VLM 决策——当任务从单个目标扩展为多房间、多模态、多目标的连续序列时，需要一个更根本的记忆重构。LagMemo（2025.10） 正是从这一需求出发，将 3DGS 提升为统一的多任务语言记忆：一次探索构建的空间-语义关联，由所有后续任务共享，无需重复建图。其设计动机源于对现有模块化导航记忆两类瓶颈的分析——基于对象的方法（如 Modular GOAT）依赖固定词汇表的上游检测器填充记忆，无法捕获开放类别目标（如"Mickey Mouse 玩偶"在探索时被忽略则永远无法检索）；基于 2D 栅格的方法将语义特征投影至平面，丢失垂直空间细节且难以跨视角保持全局一致性。
\textbf{方法信息：} 代码：\url{https://github.com/weekgoodday/LagMemo}。
\textbf{设计动机与边界：} 相对原始 3DGS 以独立高斯和图像重建误差为主的设定，LagMemo 把问题落在以下可实现的改动上：LagMemo系统包含两个核心模块：(1) 语言3DGS记忆构建——一次性探索阶段，使用前沿探索收集稀疏RGB-D和位姿序列，重建参数化3D高斯（位置、颜色、半径、不透明度），并扩展为包含码本语言特征嵌入的统一3D语言记忆。引入关键帧检索机制，按PSNR负相关权重采样历史帧进行优化，缓解大规模多房间环境中的遗忘和表面孔洞问题；(2) 记忆引导导航——对输入的多模态目标（类别、图像、文本描述），查询语言3DGS记忆预测候选目标位置，结合局部感知验证机制动态匹配和验证目标。相比2D语义地图，3DGS的连续可微性质支持实时优化和稀疏视角插值。其设计目的不是单独增加一个模块，而是用上述表示、监督或渲染约束消除原流程在目标任务中的具体失配。提出LagMemo导航系统，利用语言3DGS记忆构建统一3D语义地图，支持多模态开放词汇目标查询和多目标导航，引入目标验证机制和关键帧检索策略，在GOAT-Core基准上显著超越SOTA。适用时还应保留论文给出的边界：一次性探索阶段耗时较长；对动态环境变化的适应性有限；码本大小影响开放词汇能力。
% METHOD: 2601.13132
\paragraph{GaussExplorer（VLM 驱动的具身推理）（2026.01）}
GaussExplorer~\cite{2601_13132} 处理复杂组合语言查询。许多语言 3DGS 只支持“找到椅子”这类简单对象，遇到“靠近桌子且上面有杯子的物体”时需要推理证据和观察角度。GaussExplorer 的改变是把 VLM 放到视角选择回路中。

它先构建语义 3DGS：封闭集可用 YOLOv8，开放集可用 CLIP+SAM，然后把语义嵌入直接注册到高斯。查询时，LLM 提取“证据类别”，匹配相关语义高斯聚类，并为每个聚类选择可见性最高的初始视角。随后在邻域渲染候选新视角，让 VLM 判断哪个视角最有助于回答问题，再进行精细推理。相对 VISTA，它不只是探索覆盖，还把语言推理转成多视角观察策略。
\textbf{实现细节：} GSMem 通过反向聚合构建语言场并以几何评分（可见性、面积、不透明度）筛选视点，其视角选择完全由可计算的几何指标驱动——这在视角质量容易量化的导航场景中效果良好，但在需要理解遮挡和复杂空间关系的问答推理中可能力有不逮：预录的训练相机视角未必能提供回答问题所需的最佳证据视角。此外，GSMem 的语言场基于 CLIP 嵌入的余弦相似度检索，本质上是为类别级查询（如"椅子""桌子"）优化的静态嵌入匹配——面对组合语义查询（如"厨房里的枕头是什么颜色？"），仅靠单个高斯级别的嵌入相似度无法捕获跨对象的空间关系。GaussExplorer（2026.01） 正是针对这两重局限（固定视角约束与静态嵌入推理），将 VLM 引入视角选择循环，实现问题驱动的探索和推理。其核心洞察是：3DGS 的新视角合成能力天然充当 VLM 的"可调节相机"——不同于 GSMem 仅在已训练相机中搜索最优位姿，GaussExplorer 可在初始视角邻域内合成全新的扰动视角，由 VLM 直接评判哪个视角对回答问题最有信息量。
\textbf{方法信息：} NVIDIA 参与了这篇论文。



% METHOD: 2602.11575
\paragraph{ReaDy-Go（含移动障碍物的动态导航仿真）（2026.02）}
ReaDy-Go~\cite{2602_11575} 针对动态障碍物导航。静态 3DGS 场景可提供真实背景，但现实导航常遇到移动行人；若只在静态仿真中训练，策略很难处理动态避障。ReaDy-Go 的改变是把静态 GS 场景与可动画人体 GS 化身结合。

具体管线用 PGSR 重建目标环境的静态 GS，用 HUGS 提取参数化于 SMPL 和三平面特征的人体 GS avatar，再把人体沿 2D 轨迹插入场景并合成合理运动。结合机器人专家和人体规划器，系统生成数千个含移动人体障碍物的逼真导航场景，用于训练 RGB-only 导航策略。相对 NavGSim，它专门补动态人类障碍物；相对 EmbodiedSplat，它不只是个性化静态环境。
\textbf{实现细节：} 与 Habitat-GS 将 GS 渲染嵌入 NavMesh 导航框架不同，ReaDy-Go（2026.02） 走了另一条路线：放弃 NavMesh 和物理引擎，让动态障碍物（人类）也完全用 GS 表示，从场景到障碍物统一在 GS 表示下完成渲染和规划。这一设计选择的动机是：基于 NavMesh 的导航依赖异构表示（GS 负责视觉、NavMesh 负责碰撞），需要额外的网格提取和后端桥接；而统一 GS 表示可以实现端到端的纯视觉仿真管线——输入只需一段环境视频，即可自动生成带动态障碍物的照片级真实感导航数据集。动态 GS 仿真器。ReaDy-Go 的仿真器由两个 GS 模块组成。动态人类用 HUGS 从 NeuMan 数据集视频中提取可动画人类 GS 化身：在 SMPL 身体模型初始化的标准空间中参数化人类高斯属性（位置、SH 系数、不透明度、尺度、旋转），并学习线性混合蒙皮（LBS）权重，使人类 GS 可由 SMPL 关节参数驱动到任意姿态。
\textbf{方法信息：} 代码：\url{https://github.com/syeon-yoo/ready-go}。
\textbf{设计动机与边界：} 相对原始 3DGS 以独立高斯和图像重建误差为主的设定，ReaDy-Go 把问题落在以下可实现的改动上：ReaDy-Go包含三个核心组件：1)动态GS仿真器：用PGSR重建静态GS场景，通过HUGS提取可动画人体GS化身（参数化于SMPL身体模型+三平面特征），支持在场景中插入并沿2D轨迹生成合理人体运动；2)导航数据集生成：结合动态GS仿真器、机器人专家规划器和人体规划器，生成数千个含移动人体障碍物的逼真导航场景；3)导航策略训练：利用生成数据集训练RGB-only导航策略，同时对sim-to-real差距和动态障碍物具有鲁棒性。仅需单段视频即可构建目标环境仿真。其设计目的不是单独增加一个模块，而是用上述表示、监督或渲染约束消除原流程在目标任务中的具体失配。提出首个将静态GS场景与动态人体GS障碍物结合的真实到仿真管线，通过可动画人体GS化身和运动合成实现逼真动态导航场景生成，训练对sim-to-real差距和移动障碍物均鲁棒的导航策略。适用时还应保留论文给出的边界：依赖PGSR的高质量场景重建；人体GS化身质量受HUGS限制；2D轨迹规划可能不够自然。
% METHOD: 2603.15186
\paragraph{NavGSim（大规模导航高斯仿真器）（2026.03）}
NavGSim~\cite{2603_15186} 把 3DGS 用作大规模导航仿真器。照片级渲染只是第一步，导航还需要通行区域和碰撞检测；NavGSim 的关键改变是从重建高斯中直接提取可导航空间。

它采用分层 3DGS，把高斯组织成多层结构，导航过程中根据视角自适应选择 LoD，平衡视觉质量和速度。碰撞方面，方法用多个水平面切割高斯椭球体，得到各平面上的占用区域，从而把高斯场转换成可导航地图和碰撞表示。NavGSim-API 支持自定义场景重建、机器人配置、策略训练和评估，并兼容多 GPU。相对 LagMemo/VISTA 的语义导航，NavGSim 更偏仿真后端和策略训练环境。
\textbf{方法信息：} 代码：\url{https://github.com/2003jiahang/navgsim}。
\textbf{设计动机与边界：} 相对原始 3DGS 以独立高斯和图像重建误差为主的设定，NavGSim 把问题落在以下可实现的改动上：NavGSim基于分层3DGS框架（Hierarchical 3DGS），将3D高斯组织为多层层次结构，支持自适应LOD渲染，在导航过程中动态平衡视觉保真度和计算效率。碰撞仿真方面，引入高斯泼溅切片技术：用多个水平面切割高斯椭球体，得到的交叉区域定义物体在每个平面上的占用表示，从而将重建场景转换为可导航区域，支持高效碰撞检测和路径规划。提供NavGSim-API Python库，支持自定义场景重建（从带位姿的RGB图像）、机器人配置、策略训练和评估，兼容多GPU开发。其设计目的不是单独增加一个模块，而是用上述表示、监督或渲染约束消除原流程在目标任务中的具体失配。提出NavGSim，基于分层3D高斯泼溅的大规模导航仿真器，支持跨数百平方米场景的逼真渲染，创新性地通过高斯切片技术直接从重建高斯中提取可导航区域，并提供完整API支持多GPU开发，成功训练VLA模型并迁移到真实世界。适用时还应保留论文给出的边界：依赖高质量3DGS重建；动态物体处理有限；碰撞精度受高斯分辨率限制；计算资源需求较高。
% METHOD: 2605.21935
\paragraph{Learning to Evolve（人形导航的多模态交互场）（2026.05）}
Learning to Evolve~\cite{2605_21935} 面向人形机器人在动态环境中的导航。人形步态会带来运动模糊和抖动，语义记忆可能被低质量高斯污染；环境变化又要求地图可局部更新。该方法提出多模态交互场，把外观、空间记忆和几何验证分开。

外观场是置信度感知语义 3DGS，每个高斯携带质量置信度，用梯度幅值和不透明度识别步态损坏基元，并门控低质量语义传播。空间场从稳定渲染视图抽象拓扑记忆，通过差异检测区分运动伪变化和持久变化，只更新局部不一致区域。几何场按需重建目标网格，用于交互姿态安全验证。相对 LagMemo，它更强调非静态环境中的记忆更新和人形机器人安全约束。
\textbf{实现细节：} VR 交互与多模态记忆的表示演化。沿 Physically Embodied GS 思路扩展的 VR 交互层面，GS-Verse 与 VR-Doh 分别选取不同物理引擎实现 VR 内的变形交互：前者直接将物体网格与高斯泼溅表示集成，通过网格-高斯映射 $\Phi$ 建立高斯基元与网格顶点的关联、实现精确表面近似；后者定制物质点方法（MPM）模拟手诱导大变形，通过中轴变换近似手部几何、支持新手友好交互。沿 M3 的"记忆解耦"思路，多模态记忆沿"查询接口 $\to$ 不确定性量化 $\to$ 长期维护 $\to$ 模态扩展"四层递进：数据图方法 利用图结构允许 LMM 迭代查询大型环境的较小部分；连续语义泼溅 在 3DGS 上实现概率语义更新和光栅化，量化语义不确定性；GS-LTS 将 3DGS 从一次性重建工具转变为长期服务机器人的持续维护系统，支持动态变化室内环境的增量更新；URF-GS 提出统一射频-光学辐射场框架，为每个高斯基元扩展材质属性以支持无线信号传播仿真；MIF 提出面向人形机器人的多模态交互场，集成置信度感知语义 3DGS 与差异触发的空间记忆更新。
\textbf{实现补充：} 实现上还要核对输入、输出、参数、特征、初始化、训练损失、正则、渲染或推理接口，以及优化中的剪枝、分裂、采样、编码或解码步骤；这些抓手决定该方法到底改了 3DGS 管线的哪一环。
\textbf{技术实现补充：} VR 交互与多模态记忆的表示演化。沿 Physically Embodied GS 思路扩展的 VR 交互层面，GS-Verse 与 VR-Doh 分别选取不同物理引擎实现 VR 内的变形交互：前者直接将物体网格与高斯泼溅表示集成，通过网格-高斯映射 $\Phi$ 建立高斯基元与网格顶点的关联、实现精确表面近似；后者定制物质点方法（MPM）模拟手诱导大变形，通过中轴变换近似手部几何、支持新手友好交互。沿 M3 的"记忆解耦"思路，多模态记忆沿"查询接口 $\to$ 不确定性量化 $\to$ 长期维护 $\to$ 模态扩展"四层递进：数据图方法 利用图结构允许 LMM 迭代查询大型环境的较小部分；连续语义泼溅 在 3DGS 上实现概率语义更新和光栅化，量化语义不确定性；GS-LTS 将 3DGS 从一次性重建工具转变为长期服务机器人的持续维护系统，支持动态变化室内环境的增量更新；URF-GS 提出统一射频-光学辐射场框架，为每个高斯基元扩展材质属性以支持无线信号传播仿真；MIF 提出面向人形机器人的多模态交互场，集成置信度感知语义 3DGS 与差异触发的空间记忆更新。
\textbf{方法信息：} NVIDIA 参与了这篇论文。
\textbf{设计动机与边界：} 相对原始 3DGS 以独立高斯和图像重建误差为主的设定，Learning to Evolve 把问题落在以下可实现的改动上：MIF包含三个场：（1）外观场F\_app：置信度感知语义3DGS，每个高斯携带置信度标量Ci，通过抖动感知置信度估计识别步态损坏基元（基于归一化梯度幅值和不透明度），门控抑制低质量基元，减少运动伪影向语义传播；（2）空间场F\_spat：从稳定渲染视图抽象拓扑记忆，用于语义定位和不匹配检测，引入差异检测分数分离运动引起的假阳性变化与持久变化，仅更新局部不一致区域；（3）几何场F\_geom：基于Flow Matching的目标网格按需重建，支持交互姿态安全(IPS)验证。闭环感知-适应流水线持续稳定外观、更新记忆、验证几何。其设计目的不是单独增加一个模块，而是用上述表示、监督或渲染约束消除原流程在目标任务中的具体失配。提出MIF系统，集成置信度感知3DGS、差异触发记忆更新和交互姿态安全验证，将非静态环境重定位成功率从12%提升至94%，语义记忆占用减少91.4%适用时还应保留论文给出的边界：仅验证办公室环境；依赖RGB-D传感器；几何重建按需计算可能延迟

\section{机械臂操控}

机械臂路线通常需要对象级语义、抓取 affordance、碰撞几何和闭环仿真。3DGS 的作用可以是可编辑场景、可供性载体、机器人自模型或照片级仿真器；关键是如何把高斯从“看得见”变成“抓得到、推得动、规划得过”。
% METHOD: 2405.04378
\paragraph{Splat-MOVER（开放词汇多阶段操控）（2024.05）}
Splat-MOVER~\cite{2405_04378} 是较早把可编辑 3DGS 用于多阶段机器人操控的系统。它包含三块：ASK-Splat、SEE-Splat 和 Grasp-Splat。ASK-Splat 在预扫描阶段训练 3DGS，并把 CLIP 语义特征和抓取 affordance 特征从 2D 蒸馏到 3D，高斯因此同时带几何、语义和可操作信息。

SEE-Splat 是场景编辑模块，根据 3D 语义掩码和修复技术实时更新数字孪生，反映机器人把物体移动后的新状态。Grasp-Splat 使用 GraspNet 生成候选抓取，并用 ASK-Splat 的 affordance 分数重排序。给定自然语言多阶段任务，系统定位对象、更新场景、生成抓取并用 MoveIt 等工具规划。相对 SparseGrasp 的稀疏抓取，Splat-MOVER 更强调多阶段任务中的场景状态维护。
\textbf{实现细节：} Splat-MOVER（2024.05） 修改了 3DGS 的场景表示和渲染管线——用语义--可供性增强高斯替代标准 3D 高斯——基于 tile 的光栅化渲染机制保持不变。系统由三个模块组成：ASK-Splat 负责知识蒸馏，SEE-Splat 负责实时场景编辑，Grasp-Splat 负责可供性对齐的抓取规划。ASK-Splat：语义与可供性知识蒸馏。标准 3DGS 仅优化每个高斯的位置 $\boldsymbol{\mu}$、协方差 $\boldsymbol{\Sigma}$、不透明度 $\alpha$ 和球谐参数，缺乏语义理解和抓取知识。ASK-Splat 为此给每个高斯新增两个参数：语义特征 $f_s \in \mathbb{R}^l$ 和可供性分数 $\hat{\alpha} \in [0,1]$。然而，直接将 CLIP 的高维特征（$C=512$ 维 ViT 或 $C=1024$ 维 ResNet）赋给每个高斯会带来巨大的内存和计算开销，破坏 3DGS 的实时渲染速度。因此，ASK-Splat 引入自编码器将 CLIP 特征压缩到低维潜空间，编码器 $g_\phi^{\mathrm{enc}}: \mathbb{R}^{H_f \times W_f \times C} \to \mathbb{R}^{H_f \times W_f \times l}$ 将特征维度从 $C$ 压缩至 $l = 3$，解码器 $g_\psi^{\mathrm{dec}}$ 将其恢复。
\textbf{方法信息：} 引用很多，Semantic Scholar 约 64 次。
\textbf{设计动机与边界：} 相对原始 3DGS 以独立高斯和图像重建误差为主的设定，Splat-MOVER 把问题落在以下可实现的改动上：系统包含三个模块。ASK-Splat（Affordance-and-Semantic Knowledge）：在预扫描阶段从RGB图像训练3DGS，同时将CLIP语义特征和抓取affordance特征从2D图像蒸馏到3D场景模型中，实现几何、语义和affordance三重理解。SEE-Splat（Scene-Editing-Enabled）：基于3D语义掩码的实时场景编辑模块，使用修复（infilling）技术可视化机器人交互导致的物体运动，创建反映环境变化的数字孪生。对于多阶段任务，每完成一个阶段更新场景表示以准确反映物体新位置。Grasp-Splat：使用预训练GraspNet生成候选抓取，然后基于ASK-Splat的affordance分数重新排序获得最终抓取姿态。给定自然语言描述的多阶段任务，系统通过ASK-Splat定位和掩码物体，SEE-Splat编辑场景，Grasp-Splat生成抓取。运动规划使用MoveIt等标准工具。其设计目的不是单独增加一个模块，而是用上述表示、监督或渲染约束消除原流程在目标任务中的具体失配。提出基于可编辑3DGS的完整机器人操控栈，支持开放词汇多阶段任务。ASK-Splat实时蒸馏CLIP和抓取特征到3D场景，SEE-Splat创建数字孪生反映操控变化，Grasp-Splat生成affordance对齐的抓取。硬件实验中成功率至少提升2.4倍。适用时还应保留论文给出的边界：依赖预扫描阶段。主要针对刚性物体。实时场景编辑对复杂交互的精度有限。
% METHOD: 2409.20291
\paragraph{RL-GSBridge（Real2Sim2Real 操控强化学习）（2024.09）}
RL-GSBridge~\cite{2409_20291} 用 3DGS 构建强化学习的 Real2Sim2Real 桥梁。它先从真实桌面场景视频、相机参数和对象分割中重建仿真环境，再在仿真中训练策略并迁移回真实。

关键是 soft mesh binding GS。系统用 openMVS 生成物体 mesh，把高斯单元约束在三角面片内，通过凸组合表示位置关系，但允许颜色和其他属性更自由优化。这样既保留 GS 渲染质量，又能跟随物理仿真器中的 mesh 状态更新。仿真阶段使用 off-the-shelf 物理引擎提供动力学，GS 编辑同步视觉状态；策略用视觉感知和深度 RL 训练后部署真实抓取/放置。相对 Splatting Physical Scenes，它更偏强化学习和零样本迁移。
\textbf{方法信息：} 代码：\url{https://github.com/irmv-manipulation-group/rl-gsbridge}。
\textbf{设计动机与边界：} 相对原始 3DGS 以独立高斯和图像重建误差为主的设定，RL-GSBridge 把问题落在以下可实现的改动上：框架分为Real2Sim和Sim2Real两阶段。Real2Sim阶段：用消费级相机采集桌面场景视频，用COLMAP获取相机参数，SAM-track分割目标物体，openMVS生成网格模型。提出软网格绑定GS方法，将高斯单元约束在三角面片内，使用凸组合表示位置关系，但允许高斯属性有更大优化自由度。通过GS编辑同步物理仿真器状态。Sim2Real阶段：使用视觉感知和深度强化学习训练策略网络，直接迁移到真实世界。物理仿真采用off-the-shelf仿真器提供动态信息，GS编辑过程同步更新场景。其设计目的不是单独增加一个模块，而是用上述表示、监督或渲染约束消除原流程在目标任务中的具体失配。提出了一种基于3D Gaussian Splatting的Real2Sim2Real框架，通过软网格绑定方法增强GS渲染质量，同时保持编辑能力；结合物理仿真器实现视觉交互同步，实现零样本sim-to-real迁移。在抓取和拾取放置任务上验证了方法的有效性。适用时还应保留论文给出的边界：方法依赖于网格模型的准确性；物体类型限于桌面级操作；软绑定约束在复杂几何上可能有局限；物理仿真的精度受限于仿真器能力
% METHOD: 2412.02140
\paragraph{SparseGrasp（稀疏视角开放词汇抓取）（2024.12）}
SparseGrasp~\cite{2412_02140} 解决只有少量 RGB 视角时的开放词汇抓取。COLMAP 在 3 到 4 张图上常失败，直接 3DGS 初始化不稳；SparseGrasp 用 DUSt3R 从稀疏视角生成密集点云，再初始化 3DGS。

系统集成视觉基础模型语义特征，并用 PCA 压缩高维特征，降低存储和计算开销。场景变化时，它通过渲染当前 3DGS 并与新观测比较，快速检测变化并更新场景。抓取阶段，语言查询定位目标对象，高斯语义和几何共同提供候选，抓取规划器输出动作，支持多轮抓取。相对 3DAffordSplat 的可供性标注，它更关注稀疏输入、开放词汇和动态更新。
\textbf{方法信息：} Huawei 参与了这篇论文。
\textbf{设计动机与边界：} 相对原始 3DGS 以独立高斯和图像重建误差为主的设定，SparseGrasp 把问题落在以下可实现的改动上：SparseGrasp系统包含四个关键模块。(1)稀疏视角3DGS初始化：利用DUSt3R密集立体模型从3-4张稀疏视角图像生成密集点云作为3DGS初始化，避免COLMAP在稀疏输入下失效的问题。(2)语义感知：集成视觉基础模型的语义特征，使用PCA压缩高维特征以降低存储和计算开销。(3)高效场景更新：提出渲染-比较策略，通过渲染当前3DGS并与新观测对比快速检测场景变化，支持动态环境中的快速场景更新。(4)多轮抓取：结合语义查询和几何抓取规划，实现语言引导的多次抓取操作。整个系统在5倍速度提升下达到与密集输入方法相当的质量。其设计目的不是单独增加一个模块，而是用上述表示、监督或渲染约束消除原流程在目标任务中的具体失配。提出SparseGrasp，首个支持稀疏视角RGB图像的开放词汇机器人抓取系统。利用DUSt3R生成密集点云初始化3DGS，保持稀疏监督下的高保真。引入PCA压缩2D视觉基础模型特征，设计渲染-比较策略实现快速场景更新，支持可变环境中的多轮抓取。仅用3-4张图像即可达到SOTA性能。适用时还应保留论文给出的边界：抓取精度依赖3DGS重建质量；对极端稀疏（少于3视角）输入效果有限；物体形状复杂时抓取成功率下降。
% METHOD: 2503.05398
\paragraph{Articulated 3DGS Robot Self-Model（机器人自模型）（2025.03）}
Learning High-Fidelity Robot Self-Model with Articulated 3D Gaussian Splatting~\cite{2503_05398} 把 3DGS 用于机器人自身建模。传统 URDF 几何和纹理常较粗糙，无法真实渲染机器人外观；相对原始 3DGS 只能表示一个静态姿态，该方法的改变是把静态高斯分配到神经骨骼，再由关节角驱动变形。

流程先在静止姿态下用多视角 RGB 重建机器人 3DGS。然后提取 neural bones，每个骨骼用椭球表示，包含中心、缩放和旋转；通过聚类把高斯分给最近骨骼。运动学网络接收关节角，输出各骨骼变换矩阵，进而驱动高斯位置、旋转和尺度。训练输入是 RGB、关节角和相机参数，损失来自可微渲染图像误差；推理时输入关节角即可输出任意姿态下的高斯机器人模型，可用于运动规划、逆运动学或视觉仿真。
\textbf{实现细节：} 高保真机器人自建模（2025.03） 利用 3D 高斯表示机器人的静态形态和纹理，将 3D 高斯通过 K-means 聚类为神经椭球骨骼，每块骨骼由中心点、半径和旋转四元数定义，再通过运动学神经网络将关节角度映射为骨骼级变换。与锚点 ViserDex 的核心差异在于建模对象和变形机制：ViserDex 建模的是被操控物体的 3DGS 表示，物体变形通过施加逆变换于相机来实现（保持 3DGS 静态）；机器人自建模则建模机器人自身（而非外部物体）的 3DGS 表示，需要处理关节运动引起的非刚性变形。设计动机是机器人需要了解自身形态以进行精确的碰撞检测、运动规划和视觉位姿估计——传统方法依赖人工测量或 CAD 模型，但机器人形态会因磨损、意外损伤而变化。具体技术差异包括：（1）神经椭球骨骼表示：将 3DGS 的高斯点通过 K-means 聚类为骨骼段，每块骨骼由中心点、半径和旋转四元数参数化，运动学神经网络将关节角度映射为骨骼级 SE(3) 变换；（2）混合蒙皮权重：线性混合蒙皮（LBS）将骨骼变换传递到每个高斯，蒙皮权重由基于马氏距离的原始权重和 MLP 学习的增量权重相加后经 softmax 归一化得到——原始权重提供物理合理的初始化，增量权重补偿关节附近的非刚性变形；（3）课程学习训练：先在静态姿态训练 3DGS 并冻结高斯参数，再以课程学习方式从小范围关节运动逐步扩展到全范围运动（总计 90 万步），配合等距约束（相邻高斯间距不变）、局部刚性约束（局部高斯簇保持刚性）和循环一致性损失确保变形的物理合理性。
\textbf{方法信息：} 代码：\url{https://github.com/mer-cury-hu/selfmodel_gs}。
\textbf{设计动机与边界：} 相对原始 3DGS 以独立高斯和图像重建误差为主的设定，Learning High-Fidelity Robot Self-Model with Articulated 3D Gaussian Splatting 把问题落在以下可实现的改动上：该方法首先在静止姿态下利用多视角RGB图像重建机器人的静态3DGS模型。然后从静态模型中提取神经骨骼（neural bones），每个骨骼由椭球表示（中心点、半径缩放矩阵、旋转四元数）。通过对3D高斯进行聚类，将高斯基元分配到最近的骨骼，构建变形模型。运动学神经网络接收关节角度输入，输出各骨骼的变换矩阵，控制3D高斯的变形。训练过程使用关节角度、相机参数和多视角图像对（无需深度信息），通过可微分渲染优化3D高斯参数和运动学网络。推理时输入关节角度即可获得对应姿态下的机器人形态、运动学和纹理描述，并可从不同视角渲染图像。其设计目的不是单独增加一个模块，而是用上述表示、监督或渲染约束消除原流程在目标任务中的具体失配。提出了一种基于3D高斯溅射的机器人自建模方法，首次实现同时建模机器人形态、运动学和纹理。通过聚类3D高斯构建神经椭球骨骼，利用运动学神经网络控制变形，仅需RGB图像和关节角度即可训练，并可支持运动规划和逆运动学等下游任务。适用时还应保留论文给出的边界：需要预先知道机器人的大致运动学结构；对严重遮挡和复杂关节配置的鲁棒性有待验证；纹理建模在极端光照条件下可能受限。
% METHOD: 2504.11218
\paragraph{3DAffordSplat（高斯可供性推理）（2025.04）}
3DAffordSplat~\cite{2504_11218} 是机械臂感知路线的数据与模型代表性锚点。它构建首个面向 3DGS 的大规模 affordance 数据集，包含数万高斯实例、点云实例和多类可供性标注。原始 3DGS 没有“哪里可抓、哪里可拉、哪里可旋转”的任务语义；3DAffordSplat 把这些操作语义直接监督到高斯表示上。

流程分成数据集和模型两部分。数据侧，它为同一物体准备高斯实例、点云实例和人工可供性标签，使“高斯原语-几何点-操作语义”可以对齐；标签覆盖抓取、拉、旋转等操作类型，因此监督信号不是普通语义类别，而是面向动作的局部功能区域。模型 AffordSplatNet 的输入是 3DGS 属性和对应点云特征，核心模块是跨模态结构对齐：先把点云局部邻域和高斯局部分布映射到可比较特征空间，再用结构一致性先验训练可供性预测头，输出每个区域的操作语义热度。由于 3DGS 是连续密集分布，它能表达点云难以覆盖的细粒度接触区域。相对 SparseGrasp 的语言抓取系统，3DAffordSplat 更基础：它提供可供性数据和识别模型，为后续抓取规划提供 3D 语义热区。
\textbf{实现细节：} 改动概要：3DAffordSplat（2025.04） 修改了 3DGS 的语义表示与功能推理环节——用动态任务特定描述符替代静态语义嵌入，使每个高斯基元可根据文本查询自适应响应多种功能语义——3DGS 的几何表示和渲染管线保持不变。为此，该方法首先构建首个面向 3DGS 功能区域推理的大规模多模态数据集，包含 23,677 个高斯实例、8,354 个点云实例和 6,631 个手动标注的功能区域标签，覆盖 21 个物体类别和 18 种功能区域类型。每个物体-功能区域组合配备 15 个问答模板，支持语言引导的功能区域推理评估。方法架构：提出的 AffordSplatNet 以 3DGS 结构特征（位置 $\mu \in \mathbb{R}^{}3$、缩放 $s \in \mathbb{R}^{}3$、旋转 $r \in \mathbb{R}^{}4$）作为输入——功能区域属性主要源自物体的局部结构特征，因此模型仅处理结构特征 $\mathcal{G}_\mathrm{struct} = \{m, s, r\} \in \mathbb{R}^{10}$，外观特征（不透明度和球谐颜色）不参与推理。具体而言，利用 PointNet++ 编码多粒度几何特征，通过 RoBERTa 语言模型提取文本查询的 $\langle\mathrm{Aff}\rangle$ token 嵌入，该 token 同时编码了问题理解和掩码生成的中间语义表示。跨模态融合模块在空间和通道两个层面执行交叉注意力——空间注意力以语言特征为查询、几何特征为键值，通道注意力通过拼接和残差连接融合全局语言上下文与局部几何细节。粒度自适应选择机制通过可学习门控权重 $W_\mathrm{gate}$ 在三个粒度级别上竞争性地分配重要性，最终通过逆距离加权插值上采样至原始高斯密度。解码器基于 Transformer 生成动态卷积核，将多模态特征解码为逐高斯的功能区域掩码——这意味着卷积核不是固定的，而是根据当前文本查询动态生成，从而实现同一高斯基元在不同功能上下文中产生不同的功能区域预测。
\textbf{方法信息：} 代码：\url{https://github.com/hcplab-sysu/3daffordsplat}；GitHub 关注度较高，约 79 stars。



% METHOD: 2506.19842
\paragraph{ManiGaussian++（双臂层次高斯世界模型）（2025.06）}
ManiGaussian++~\cite{2506_19842} 把高斯世界模型用于双臂操作。双臂任务的难点是稳定臂和操作臂会造成不同类型的场景动态，普通单体视觉表示难以建模多体时空耦合。方法提出 Leader-Follower 层次高斯世界模型。

输入 RGB-D 经过体素化和稀疏卷积编码成视觉表示，再生成任务导向的高斯辐射场。Leader 分支预测稳定臂运动引起的高斯形变，Follower 分支预测操作臂动作的物理后果；每个高斯粒子携带任务属性。训练目标结合重建损失、未来动态预测损失和策略损失，使中间高斯表示同时服务场景预测和动作决策。相对 GWM 的通用动作条件未来预测，ManiGaussian++ 更针对双臂协同结构。
\textbf{方法信息：} 代码：\url{https://github.com/april-yz/manigaussian_bimanual}。
\textbf{设计动机与边界：} 相对原始 3DGS 以独立高斯和图像重建误差为主的设定，ManiGaussian++ 把问题落在以下可实现的改动上：核心创新：(1)面向任务的高斯泼溅：从视觉特征生成任务导向的高斯辐射场，区分稳定臂和操作臂以建模多体时空动态。每个高斯粒子被分配特定任务属性。(2)层次化高斯世界模型：Leader-Follower架构——Leader预测稳定臂运动引起的高斯泼溅变形，Follower生成操作臂运动的物理后果。通过未来场景预测将多体时空动态编码到中间视觉表示中。(3)学习目标：结合重建损失、动态预测损失和策略损失进行端到端训练。从RGB-D输入转换为体素空间，通过稀疏卷积网络编码为视觉表示。其设计目的不是单独增加一个模块，而是用上述表示、监督或渲染约束消除原流程在目标任务中的具体失配。扩展ManiGaussian至双臂场景，提出层次化高斯世界模型（Leader-Follower架构），分别建模稳定臂和操作臂的多体时空动态，在10个仿真任务中提升20.2%，9个真实任务中平均成功率60%。适用时还应保留论文给出的边界：依赖RGB-D输入和精确标准；对未见过的物体泛化能力有限；层次化架构增加了训练复杂度。
% METHOD: 2510.20813
\paragraph{GSWorld（闭环照片级操控仿真器）（2025.10）}
GSWorld~\cite{2510_20813} 是机械臂仿真体系方法。它提出 GSDF 资产格式，把 Gaussian-on-Mesh 表示、机器人 URDF、物体属性和碰撞网格放在同一资产中。这样策略训练、DAgger 数据收集和基准评估能使用同一套照片级数字孪生。

Real-to-Sim 流程从多视角捕获构建度量准确场景，用 ArUco 标记确定绝对尺度，避免 COLMAP 尺度模糊；用 ICP 对齐机器人 URDF 到场景，并用 2DGS 提升几何精度。Sim-to-Real 流程保持仿真和真实机器人的控制/观察空间一致，使训练策略可直接部署硬件。数据库包含多种机器人和四十余个物体。相对 GaussGym 的通用向量化仿真，GSWorld 更聚焦机械臂闭环操控资产和评测。
\textbf{实现细节：} GSWorld（2025.10） 提出 GSDF（Gaussian Scene Description File）格式，将高斯-网格表示与机器人 URDF 融合，支持零样本仿真到真实操控策略学习、自动 DAgger 数据收集和可复现策略基准评测。DISCOVERSE（2025.07） 提出首个统一、模块化、开源的 3DGS 仿真框架，通过全面 Real2Sim 管线合成超逼真几何和外观。
\textbf{实现补充：} 实现上还要核对输入、输出、参数、特征、初始化、训练损失、正则、渲染或推理接口，以及优化中的剪枝、分裂、采样、编码或解码步骤；这些抓手决定该方法到底改了 3DGS 管线的哪一环。
\textbf{方法信息：} NVIDIA 参与了这篇论文。
\textbf{设计动机与边界：} 相对原始 3DGS 以独立高斯和图像重建误差为主的设定，GSWorld 把问题落在以下可实现的改动上：GSWorld框架包含：(1)GSDF资产格式：将Gaussian-on-Mesh表示与机器人URDF和物体属性融合，使用ArUco标记设定绝对尺度（避免COLMAP尺度模糊），通过ICP对齐机器人URDF到场景，附加碰撞网格和材质属性；(2)Real-to-Sim流程：从多视角捕获构建度量精确的数字孪生，使用2DGS提高几何精度；(3)Sim-to-Real流程：策略在sim中训练后直接部署到硬件（控制和观察空间匹配机器人原生API）。支持多种应用：零样本sim-to-real像素到动作策略学习、DAgger自动数据收集、可复现基准测试、虚拟遥操作数据收集、零样本sim-to-real视觉强化学习。数据库包含3种机器人实体（单臂和双臂）和40+物体。其设计目的不是单独增加一个模块，而是用上述表示、监督或渲染约束消除原流程在目标任务中的具体失配。提出GSWorld闭环照片级仿真器和GSDF资产格式，集成3DGS渲染与物理引擎，支持零样本sim-to-real策略迁移、DAgger数据收集和可复现基准测试，构建了3种机器人+40+物体的GSDF数据库。适用时还应保留论文给出的边界：GSDF重建质量依赖输入图像质量；物理引擎与3DGS的对齐仍有偏差；大规模场景重建耗时
% METHOD: 2604.11138
\paragraph{ViserDex（高斯空间域随机化的灵巧手）（2026.04）}
ViserDex~\cite{2604_11138} 处理单目 RGB 灵巧手物体重定向。Sim2Real 的视觉差距常来自渲染不真实和域随机化太粗糙；ViserDex 的改变是在 3DGS 表示空间中做物理一致的域随机化，再渲染训练物体位姿估计器。

管线分三段。先在仿真中用 PPO 和 asymmetric actor-critic 训练有完整状态的教师策略。再把教师蒸馏成循环学生策略，让它从噪声观测推断潜在状态。最后训练视觉位姿估计器：在渲染前对物体 3D 高斯施加颜色、纹理、光照和背景等扰动，生成照片级但多样的训练图像。部署时，RGB 位姿估计结果输入学生策略，完成多指手重定向。相对 RL-GSBridge 的桌面抓取，它面向更难的 in-hand dexterous manipulation。
\textbf{实现细节：} ViserDex（2026.04） 修改了 3DGS 的渲染管线数据增强环节——用预光栅化增强（pre-rasterization augmentation）替代传统后处理图像增强——3D 高斯的位置、协方差、不透明度保持不变，仅修改球谐系数以生成视觉多样的训练图像。该方法解决的核心问题是灵巧操控中的仿真到真实迁移：多指手操控物体时，物体在手指间快速运动、频繁被遮挡且伴随运动模糊，导致直接从 RGB 图像估计物体位姿极为困难。传统生成训练数据的策略面临两难——后处理图像增强（如颜色抖动、亮度变化）计算便宜但仅施加全局 2D 变换，无法捕捉 3D 空间中光照和材质的结构化变化；而基于光线追踪的域随机化虽然可以精确模拟光照变化，但计算开销过大，难以在仿真循环中实时运行。ViserDex 的核心洞察在于：3DGS 的显式高斯属性——特别是球谐（SH）系数——本身就是一个天然的数据增强操作界面。每个 3D 高斯的颜色由球谐系数决定，其中零阶系数 $\mathrm{SH0}$ 编码与视角无关的漫反射基色，高阶系数 $\mathrm{SHN}$ 编码高频镜面效果。
\textbf{技术实现补充：} ViserDex（2026.04） 修改了 3DGS 的渲染管线数据增强环节——用预光栅化增强（pre-rasterization augmentation）替代传统后处理图像增强——3D 高斯的位置、协方差、不透明度保持不变，仅修改球谐系数以生成视觉多样的训练图像。该方法解决的核心问题是灵巧操控中的仿真到真实迁移：多指手操控物体时，物体在手指间快速运动、频繁被遮挡且伴随运动模糊，导致直接从 RGB 图像估计物体位姿极为困难。传统生成训练数据的策略面临两难——后处理图像增强（如颜色抖动、亮度变化）计算便宜但仅施加全局 2D 变换，无法捕捉 3D 空间中光照和材质的结构化变化；而基于光线追踪的域随机化虽然可以精确模拟光照变化，但计算开销过大，难以在仿真循环中实时运行。ViserDex 的核心洞察在于：3DGS 的显式高斯属性——特别是球谐（SH）系数——本身就是一个天然的数据增强操作界面。每个 3D 高斯的颜色由球谐系数决定，其中零阶系数 $\mathrm{SH0}$ 编码与视角无关的漫反射基色，高阶系数 $\mathrm{SHN}$ 编码高频镜面效果。论文提出的预光栅化增强直接在高斯属性空间操作，在光栅化之前对 SH 系数施加结构化随机扰动，既保持了物理一致性（相邻高斯协同变化），又避免了全场景光线追踪的计算开销。具体而言，论文设计了四层增强策略，从局部到全局逐层施加：（1）随机噪声层：对每个高斯独立施加随机扰动，模拟高频传感器噪声和像素级伪影；（2）空间聚类层：将高斯按空间位置用 k-means 聚为 $K_s=64$ 个空间簇，对每个簇整体扰动，模拟阴影、磨损等局部一致的外观变化；（3）颜色聚类层：将高斯按 $\mathrm{SH0}$ 基色聚为 $K_c=32$ 个颜色簇，对同一材质区域的漫反射属性施加统一偏移，模拟反照率变化或反射率差异；（4）全局偏移层：将整个高斯场景视为单一簇，施加全局亮度和色温偏移，模拟环境光照变化和相机曝光差异。每层增强以概率 $p_\mathrm{aug}$ 激活，仅扰动比例为 $p_\mathrm{cluster}$ 的簇，扰动幅度在 $[\eta_\mathrm{min}, \eta_\mathrm{max}]$ 范围内均匀采样，多层扰动可叠加产生复合效果。渲染与遮挡合成。在灵巧操控中，相机固定而物体被手指操控产生运动。
\textbf{方法信息：} NVIDIA 参与了这篇论文。
\textbf{设计动机与边界：} 相对原始 3DGS 以独立高斯和图像重建误差为主的设定，ViserDex 把问题落在以下可实现的改动上：模块化流程分为三个阶段：(1)教师训练：在仿真中使用完整状态信息的强化学习(PPO)训练教师策略，使用Asymmetric Actor-Critic框架，仅优化操控技能；(2)学生蒸馏：将教师蒸馏为循环学生策略，学习从噪声观测推断潜在系统状态；(3)视觉位姿估计器训练：核心创新是在3DGS表示空间中进行域随机化——对物体的3D高斯应用物理一致的颜色扰动、纹理变换、光照变化和背景替换等增强，然后通过3DGS渲染生成照片级真实的训练图像。部署时，从RGB图像预测物体位姿，输入循环学生策略产生目标条件动作。其设计目的不是单独增加一个模块，而是用上述表示、监督或渲染约束消除原流程在目标任务中的具体失配。提出ViserDex，首个在3DGS表示空间中执行域随机化的灵巧操控Sim2Real框架。通过在渲染前对3D高斯应用物理一致的数据增强，生成照片级真实且多样化的视觉数据训练物体位姿估计器。结合课程式强化学习和教师-学生蒸馏训练操控策略。两个模型均可在消费级硬件上独立训练，无需大规模计算集群。在物理多指手上验证5种物体的鲁棒重定向。适用时还应保留论文给出的边界：假设刚体物体，仅验证了5种物体，对非刚体和透明物体的适用性未知，位姿估计精度受3DGS表示质量限制。

\section{具身感知与交互}

这一节的方法把 3DGS 从静态地图提升为可预测、可校正、可交互或跨模态的世界模型。关键增量是给高斯附加物理粒子、语义不确定性、无线辐射属性、LLM 交互能力或 VR 物理绑定。
% METHOD: 2406.10788
\paragraph{Physically Embodied Gaussian Splatting（实时可校正物理世界模型）（2024.06）}
Physically Embodied Gaussian Splatting~\cite{2406_10788} 是预测-校正世界模型的早期锚点。它提出 dual Gaussian-particle representation：粒子系统承载几何和物理状态，每个粒子附着一个 3D Gaussian 用于视觉渲染。

运行时，粒子物理系统先预测未来状态；高斯 splatting 从任意视角渲染预测图像；预测图像和观测图像之间的差异被转换成视觉力，在物理约束下校正粒子位置。这样，系统既能做物理预测，又能通过视觉反馈保持与真实环境同步。仅用少数相机即可进行 2D/3D 跟踪和光度重建验证，并达到实时频率。相对 GaussTwin，它更早提出视觉力校正思想；GaussTwin 后续把刚体和 Cosserat 杆建模得更具体。
\textbf{实现细节：} Physically Embodied Gaussian Splatting（2024.06） 修改了 3DGS 的场景表示与状态更新机制——用"粒子-高斯"双表示替代纯高斯表示，通过物理引擎驱动的预测-校正闭环替代直接梯度优化——高斯基元的可微泼溅渲染流程保持不变。该方法的核心创新在于引入"高斯-粒子"（Gaussian-Particle）双表示，用一个刚体变换（bond，即记录高斯基元相对于父粒子的局部坐标偏移与朝向）将每个高斯基元挂载到其父粒子上，从而在物理仿真和可微渲染之间建立桥梁。直观地说，每个高斯基元就像用刚性连杆悬挂在粒子上的灯笼——灯笼负责向观察者展示物体的外观，连杆则保证灯笼始终跟随粒子的物理运动轨迹。具体而言，每个场景包含两类实体：$N$ 个粒子负责捕获物体的几何与物理状态（位置、速度、质量等），由 Position-Based Dynamics（PBD，一种以位置为直接求解变量的实时物理仿真方法，通过迭代满足约束来模拟刚体碰撞、变形等行为）物理引擎驱动；$M$ 个高斯基元负责从任意视角渲染图像，其位置随父粒子刚性移动。一个典型桌面场景约有 1000 个粒子和 10000 个高斯基元。系统的运行分为预测和校正两个阶段，形成"物理仿真 $\to$ 渲染比较 $\to$ 视觉力反馈"的闭环。在预测阶段，PBD 物理引擎以 30\,Hz 对粒子执行前向仿真：首先根据机器人正运动学设定机器人粒子的位置，然后通过 Jacobi 迭代求解器同时处理三类约束——地面约束（防止粒子穿透地面平面）、碰撞约束（防止粒子间相互穿透）和形状匹配约束（保持刚体或可变形体的结构），每步物理仿真约耗时 5\,ms。随后，挂载的高斯基元随父粒子刚性移动至预测位置。
\textbf{实现补充：} 实现上还要核对输入、输出、参数、特征、初始化、训练损失、正则、渲染或推理接口，以及优化中的剪枝、分裂、采样、编码或解码步骤；这些抓手决定该方法到底改了 3DGS 管线的哪一环。
\textbf{方法信息：} 引用很多，Semantic Scholar 约 52 次。
\textbf{设计动机与边界：} 相对原始 3DGS 以独立高斯和图像重建误差为主的设定，Physically Embodied Gaussian Splatting 把问题落在以下可实现的改动上：该方法的核心是Dual Gaussian-Particle Representation。粒子系统负责捕捉场景中物体的几何形状，并与基于粒子的物理系统配合，预测物理上合理的未来状态。每个粒子附着一个3D高斯，通过splatting过程从任意视角渲染图像以捕捉视觉状态。关键创新在于利用预测图像与观测图像之间的差异生成「视觉力」，这些力在遵守已知物理约束的条件下校正粒子位置。这种统一表示使得系统既能在物理层面上推理现在和未来，又能与现实保持同步。系统使用3个相机即可运行，频率达到30Hz的实时性能。整体框架将预测性物理建模与连续的视觉衍生校正整合在一起，实现了物理仿真与视觉感知的闭环。其设计目的不是单独增加一个模块，而是用上述表示、监督或渲染约束消除原流程在目标任务中的具体失配。提出双高斯-粒子表示法，首次将物理预测仿真与视觉在线校正统一到一个框架中。系统能以30Hz实时运行，仅需3个相机即可实现2D和3D跟踪以及光度重建质量验证，为机器人在动态环境中提供了同时预测和校正的能力。适用时还应保留论文给出的边界：目前仅在特定场景中验证，对复杂多物体交互的处理能力有待进一步探索；物理仿真的精度受限于粒子系统的表达能力。
% METHOD: 2409.16019
\paragraph{AIR-Embodied（具身 LLM 主动重建）（2024.09）}
AIR-Embodied~\cite{2409_16019} 将 3DGS 与具身大模型结合，用于主动重建。普通主动建图通常只选择新视角；遮挡严重时，机器人还需要移动物体才能看到被遮挡区域。AIR-Embodied 的改变是让 MLLM 读取重建状态并规划视角和交互动作。

系统从像素级重建质量生成多模态提示，利用 MLLM 判断哪些区域质量差或被遮挡。规划阶段联合优化视角选择和动作选择，目标包含重建损失和视角/动作成本；动作可以是推动物体以暴露遮挡面。闭环验证阶段比较计划与实际执行结果，动态修正。底层使用 PGSR 变体的扁平化 3DGS 表示，并用沿深度透射率分布的不确定性作为质量代理。相对 VISTA 的任务相关探索，AIR-Embodied 更强调物理交互帮助重建。
\textbf{实现细节：} 与 SplatR 在已完成重建的静态场景中进行物体重排的任务设定不同，AIR-Embodied（2024.09） 将 3DGS 的角色从被动场景表示前置为主动重建过程的优化目标——其核心问题是：如何利用多模态大语言模型（MLLM）的高层推理能力，规划视点选择与交互动作以最大化 3DGS 的重建质量？该问题的瓶颈在于，传统 Next Best View（NBV，下一最优视角）方法（如基于不确定性的和基于 Fisher 信息的方法）仅从球面候选集中选择视点，依赖低层不确定性信号而缺乏对重建状态的高层理解，且无法处理物体自身遮挡（如物体底部）——即使选定最优视点，接触面仍不可见。AIR-Embodied 的设计动机正是突破这两层限制：(1) 用 MLLM 将像素级不确定性提升为语义级理解，实现自由空间中的视点规划而非限于球面候选集；(2) 引入机械臂交互操作暴露遮挡区域，使重建不再受限于被动观测。三阶段主动重建框架。AIR-Embodied 采用 PGSR（平面正则化 3DGS 变体）作为底层面表示，在其上构建三阶段闭环管线。阶段一：重建状态理解。
\textbf{实现补充：} 实现上还要核对输入数据、输出指标、基线配置、评测协议、系统接口和适用边界；若调用 3DGS 渲染或重建结果，也应说明相机参数、训练设置、优化流程和指标计算如何对齐。
\textbf{方法信息：} 代码：\url{https://github.com/qzh-00/air-embodied}。



% METHOD: 2411.02547
\paragraph{Continuous Semantic Splatting（语义不确定性建图）（2024.11）}
Continuous Semantic Splatting~\cite{2411_02547} 关注安全关键机器人中的不确定性。标准 3DGS 可以把语义标签渲染出来，但通常只给点估计，不给置信度。CSS 的核心洞察是 3DGS 的 alpha blending 贡献函数满足贝叶斯核推断约束，因此可直接用于概率语义更新。

每个 3D 椭球体持有共轭先验分布，2D 语义分割网络的预测作为观测对其进行贝叶斯更新，得到语义期望和方差。概率光栅化再把这些先验分布投到像素空间，输出带不确定性的语义分割。方法不需要额外训练，也避免体素离散化误差。相对 LagMemo/VISTA 的语言语义地图，CSS 的重点是可量化风险和置信度。
\textbf{实现细节：} 连续语义泼溅（2024.11） 在 3DGS 上实现概率语义更新和光栅化。与 M3 通过注意力机制蒸馏基础模型特征不同，该方法改变了 3DGS 的语义编码和渲染环节——利用 3DGS 的 $\alpha$ 混合权重作为贝叶斯核推理（BKI）的核函数，直接在 3D 椭球体上执行概率语义更新，而非通过额外训练学习特征渲染。设计动机在于：传统语义建图方法（如体素化 BKI）将语义信息离散化至固定分辨率体素，存在精度和内存的矛盾，而 3DGS 的椭球体天然提供连续、自适应分辨率的场景覆盖，可以作为 BKI 核函数的直接载体。具体而言，该方法在每个椭球体上维护 Dirichlet 共轭先验分布，通过 $\alpha_i^n = \alpha_i^{n,0} + \sum_j w(x_j, x_n) y_j^c$ 累积语义观测浓度参数，其中 $w(x_j, x_n)$ 即 3DGS 的 $\alpha$ 混合权重，在满足 BKI 核函数约束（$0 \leq w \leq 1$、$\sum w = 1$）的同时无需额外训练。概率光栅化通过 $E[\pi_i] = \sum_n w(x_i, x_n) E[\pi_n]$ 和方差传播 $V[\pi_i] = \sum_n w(x_i, x_n)^{}2 V[\pi_n]$ 返回可量化不确定性的逐像素分割预测，将连续语义建图从体素扩展至椭球体。
\textbf{方法信息：} Amazon 参与了这篇论文。
\textbf{设计动机与边界：} 相对原始 3DGS 以独立高斯和图像重建误差为主的设定，Modeling Uncertainty in 3D Gaussian Splatting through Continuous Semantic Splatting 把问题落在以下可实现的改动上：方法基于一个核心洞察：3DGS中的alpha blending贡献函数满足贝叶斯核推断（BKI）的约束条件，因此3DGS学习到的核可直接用于概率更新。方法分三步：(1) 概率语义更新——对每个3D椭球体使用共轭先验进行贝叶斯更新，通过语义分割网络的预测作为观测，获得每个椭球体的语义期望和方差；(2) 概率光栅化——将3D共轭先验分布光栅化到像素空间，保持可量化的不确定性；(3) 无需额外训练——直接利用已训练的3DGS模型，不需要学习渲染特征。与基于体素的方法相比，避免了离散化误差。其设计目的不是单独增加一个模块，而是用上述表示、监督或渲染约束消除原流程在目标任务中的具体失配。提出CSS方法，首次将连续语义建图从体素扩展到3DGS椭球体表示。核心洞察是3DGS学习到的核函数满足BKI约束，可直接用于概率贝叶斯更新。设计概率光栅化方法，返回带有可量化不确定性的像素级语义分割预测，为安全关键机器人应用提供置信度信息。适用时还应保留论文给出的边界：依赖预训练语义分割网络的准确性；在稀疏视图区域不确定性估计可能不够准确；3DGS重建质量影响语义建图精度。
% METHOD: 2508.17600
\paragraph{GWM（动作条件 3D 高斯世界模型）（2025.08）}
GWM~\cite{2508_17600} 把 3DGS 用作机器人操作的未来状态预测器。它不是只重建当前场景，而是学习在机器人动作作用下高斯基元如何传播到未来。输入可来自单/双视角图像，先由 Splatt3R 前馈重建 3D 高斯，再进入高斯 VAE。

3D 高斯 VAE 将数量不一的高斯编码为固定长度潜嵌入，使用 Transformer cross-attention 编码器和镜像解码器。潜在 DiT 在这个潜空间中学习动作条件扩散，预测未来高斯状态。该模型有两种用途：离线模仿学习中，自监督未来预测增强视觉表示；在线模型强化学习中，GWM 作为神经仿真器生成训练数据。相对 ManiGaussian++ 的双臂特化，GWM 更像通用 3D 世界模型。
\textbf{方法信息：} 引用很多，Semantic Scholar 约 38 次。
\textbf{设计动机与边界：} 相对原始 3DGS 以独立高斯和图像重建误差为主的设定，GWM 把问题落在以下可实现的改动上：方法包含四个核心组件：(1)前馈3DGS编码：使用Splatt3R从单/双视角图像快速重建3D高斯表示，避免耗时的逐场景优化；(2)3D高斯VAE：将数量不一的3D高斯编码为固定长度N的潜嵌入，使用Transformer交叉注意力编码器和镜像Transformer解码器；(3)潜在DiT：在潜空间中使用条件扩散模型学习动作条件化的场景动态，预测未来高斯状态；(4)策略学习：支持两种模式——离线模仿学习（自监督未来预测增强视觉表示）和在线模型强化学习（作为神经仿真器生成训练数据）。其设计目的不是单独增加一个模块，而是用上述表示、监督或渲染约束消除原流程在目标任务中的具体失配。提出首个3D高斯世界模型GWM用于机器人操作，通过推断高斯基元在机器人动作作用下的传播来重建未来状态。核心是结合潜在扩散Transformer(DiT)和3D变分自编码器，实现细粒度场景级未来状态重建。GWM既可作为自监督未来预测训练增强模仿学习的视觉表示，也可作为神经仿真器支持基于模型的强化学习。展示了3D世界模型的初始数据扩展潜力。适用时还应保留论文给出的边界：前馈3DGS重建质量受限；对复杂多步操作的长期预测精度可能下降；训练数据需求较大。
% METHOD: 2510.11878
\paragraph{GS-Verse（VR 中的物理感知高斯交互）（2025.10）}
GS-Verse~\cite{2510_11878} 面向 VR 中的交互式高斯资产。过去一些 VR-GS 方法用四面体笼或简化代理做物理交互，表面近似粗，拉伸或扭曲不真实。GS-Verse 直接把高斯绑定到物体表面 mesh，并让 mesh 作为物理模拟的主要表示。

当用户在 VR 中拉伸、扭曲或摇晃物体时，物理引擎更新表面 mesh，高斯点随绑定关系发生形变和运动，渲染仍由 3DGS 保持细节。该设计能利用已有 3D mesh 资产，且不绑定特定物理引擎，可在 Unity 等平台部署。相对本章机器人操控系统，它偏人机交互和沉浸式资产操作，但技术上展示了高斯与物理 mesh 的可交互绑定。
\textbf{实现细节：} GS-Verse（2025.10） 直接将物体网格与高斯泼溅表示集成，通过网格-高斯映射 $\Phi$ 建立高斯基元与网格顶点之间的关联，在 VR 环境中实现更精确的表面近似和高度真实的变形交互。18 人用户研究证实该方法在物理感知拉伸操控上统计显著优于 VR-GS 基线，同时在系统可用性（SUS）和认知负荷（NASA-TLX）方面表现优异。VR-Doh（2024.12） 定制物质点方法（MPM）模拟手诱导大变形，通过中轴变换（MAT）近似手部几何实现高效碰撞处理，结合解耦外观与物理表示和均匀高斯体积正则化技术增强 3DGS 无缝渲染，支持新手用户直观的接触式和手势式 3D 建模体验。
\textbf{实现补充：} 实现上还要核对输入、输出、参数、特征、初始化、训练损失、正则、渲染或推理接口，以及优化中的剪枝、分裂、采样、编码或解码步骤；这些抓手决定该方法到底改了 3DGS 管线的哪一环。
\textbf{方法信息：} 代码：\url{https://github.com/anastasiya999/gs-verse}。
\textbf{设计动机与边界：} 相对原始 3DGS 以独立高斯和图像重建误差为主的设定，GS-Verse 把问题落在以下可实现的改动上：GS-Verse的核心创新在于直接使用物体表面网格作为物理模拟的主要表示，而非之前VR-GS等方法使用的简化四面体笼。该方法将高斯溅射点绑定到物体表面网格上，通过网格的物理模拟驱动高斯点的形变和运动。这使得：(1)更精确的表面近似，实现高度真实的形变；(2)利用现有3D网格资产，简化开发工作流；(3)物理引擎无关设计，支持多种物理引擎部署；(4)支持拉伸、扭曲、摇晃等多种物理感知操作。系统在Unity等VR开发平台上实现，提供沉浸式的3D内容交互体验。其设计目的不是单独增加一个模块，而是用上述表示、监督或渲染约束消除原流程在目标任务中的具体失配。提出了GS-Verse方法，直接将物体的网格与高斯溅射表示集成，实现更精确的表面近似和高度真实的形变与交互。与现有依赖四面体笼等简化几何代理的方法不同，GS-Verse利用物体表面网格作为物理模拟的主要表示，支持物理引擎无关部署。用户研究（18名参与者）显示在物理感知拉伸操作上显著优于SOTA。适用时还应保留论文给出的边界：依赖高质量物体网格资产；对复杂拓扑物体的网格绑定可能具有挑战性；VR设备的计算能力限制场景规模。
% METHOD: 2601.19216
\paragraph{URF-GS（视觉与无线统一辐射场）（2026.01）}
Bridging Visual and Wireless Sensing via a Unified Radiation Field~\cite{2601_19216} 是跨模态具身感知方法。机器人和数字孪生不只需要光学地图，也可能需要无线电地图；URF-GS 的改变是用共享高斯基元同时表示视觉几何和无线传播属性。

光学域中，方法用单目模型估计深度和法线作为伪监督，提高 3DGS 几何和多视图一致性。无线域中，它把角功率或空间频谱通过全景投影转成针孔图像形式，再用物理信息逆向渲染模拟散射和传播。训练时先固定几何，用频谱监督优化高斯材质属性；推理时给定任意收发位置即可预测无线信号，无需重训。相对机器人视觉导航地图，它服务的是视觉-无线一体化环境感知。
\textbf{实现细节：} URF-GS（2026.01） 提出统一射频-光学辐射场框架。与 M3 在视觉语言模态上构建多模态记忆不同，URF-GS 改变了 3DGS 的属性表示和渲染环节——为每个高斯基元扩展材质属性（散射反照率、金属度、粗糙度）和法向量，使同一组高斯同时支持光学渲染和射频信号传播仿真。设计动机在于：机器人在复杂环境中导航和规划时不仅依赖视觉信息，还需要预测无线信号的空间分布（如 Wi-Fi 接入点部署优化），而传统方法将视觉和无线建模割裂处理，无法共享场景几何信息。URF-GS 通过两阶段管线解决这一问题：第一阶段利用单目深度和法向量先验正则化 3DGS 优化以获取精确几何；第二阶段基于物理的逆向渲染将射频信号传播分解为表面交互（散射+反射，遵循 Cook-Torrance BRDF 模型）和自由空间传输（自由空间路径损耗模型），从多模态观测联合优化材质属性和辐射模式。
\textbf{技术实现补充：} URF-GS（2026.01） 提出统一射频-光学辐射场框架。与 M3 在视觉语言模态上构建多模态记忆不同，URF-GS 改变了 3DGS 的属性表示和渲染环节——为每个高斯基元扩展材质属性（散射反照率、金属度、粗糙度）和法向量，使同一组高斯同时支持光学渲染和射频信号传播仿真。设计动机在于：机器人在复杂环境中导航和规划时不仅依赖视觉信息，还需要预测无线信号的空间分布（如 Wi-Fi 接入点部署优化），而传统方法将视觉和无线建模割裂处理，无法共享场景几何信息。URF-GS 通过两阶段管线解决这一问题：第一阶段利用单目深度和法向量先验正则化 3DGS 优化以获取精确几何；第二阶段基于物理的逆向渲染将射频信号传播分解为表面交互（散射+反射，遵循 Cook-Torrance BRDF 模型）和自由空间传输（自由空间路径损耗模型），从多模态观测联合优化材质属性和辐射模式。
\textbf{方法信息：} Meta 参与了这篇论文。
\textbf{设计动机与边界：} 相对原始 3DGS 以独立高斯和图像重建误差为主的设定，Bridging Visual and Wireless Sensing via a Unified Radiation Field for 3D Radio Map Constr 把问题落在以下可实现的改动上：框架核心创新：(1)共享高斯基元统一表示光学和无线域。光学域使用单目模型估计深度和法线作为伪GT监督几何属性，增强3DGS的深度恢复和多视图一致性。(2)无线域使用全景投影将角功率/空间频谱转为针孔相机图像。物理信息逆向渲染循环分解散射现象模拟波传播。(3)训练时频谱监督高斯的材质属性（几何冻结），推理时仅需位置输入即可重建光学图像和空间频谱。支持任意Tx-Rx位置配置，无需重训练。其设计目的不是单独增加一个模块，而是用上述表示、监督或渲染约束消除原流程在目标任务中的具体失配。提出URF-GS统一视觉-无线辐射场框架，基于3DGS和逆向渲染构建3D无线电地图。通过融合跨模态观测恢复场景几何和材质属性，在任意收发配置下预测无线电信号无需重训练。相比NeRF方法，空间频谱精度提升24.7%，采样效率提升10倍。适用时还应保留论文给出的边界：依赖场景几何的准确恢复；复杂电磁环境（如强多径效应）的建模精度有待验证；目前假设各向同性材质属性。

\section{主动建图、具身学习与数据生成}

这一节收纳更靠近具身智能数据闭环的方法：多机器人下一最佳视角、VLA 微调、人形机器人 synthetic interaction、视觉语言导航统一评测，以及从 RGB-D 视频恢复可交互铰接物体。它们共同把 3DGS 当作可生成、可评估或可规划的环境介质。
% METHOD: 2402.07181
\paragraph{3D Gaussian as a New Era（机器人、导航与具身任务）（2024.02）}
3D Gaussian as a New Era~\cite{2402_07181} 关注的问题是：把 3DGS 用于机器人、导航或具身智能 3D Gaussian as a New Era 是宽召回补全中发现的论文，主要根据标题、摘要和 arXiv 类别判定其与 3DGS 的关系。
相对原始 3DGS，3D Gaussian as a New Era 的变化不是只替换损失函数，而是把表示、初始化、约束或网络预测路径重新组织到任务流程中：该工作被召回是因为论文文本包含 Gaussian/3D Gaussian 相关表述；当前分类先确定其任务域和与 3DGS 的相关性，后续证据流水线再补代码、企业和引用信息。
\textbf{实现细节：} 具体实现可以按输入、核心模块和渲染/优化路径来读。输入端首先取得任务所需的图像、相机、文本、时间序列、语义或几何线索，并把这些线索转换成可以约束高斯中心、尺度、旋转、不透明度、颜色或附加特征的条件。核心结构中，方法使用的主要机制是：从摘要层面看，3D Gaussian as a New Era 围绕“把 3DGS 用于机器人、导航或具身智能”展开。若其正文或代码后续可读，应进一步核对它具体改动的是表示、初始化、渲染、优化、损失、动态建模、压缩或应用接口中的哪一环。 训练或推理时，这些模块通常先给出初始化或中间特征，再通过可微 Gaussian rasterization、重建损失、感知/语义约束、几何正则或任务损失回传到高斯参数；如果方法包含前馈网络，则网络输出的是高斯属性、变形场、压缩码、语义特征或控制变量，而不是只输出最终 2D 图像。渲染路径仍以显式高斯投影、深度排序和 alpha 合成为基础，只是在采样、属性预测、局部编辑、动态变形、压缩编码或下游任务接口处加入额外约束。这样写入本章时，它与前后方法的关系是：相对原始 3DGS 和本章已有锚点方法，它可以理解为 3DGS extends\_or\_uses 宽召回分类阶段的关系判断，后续需要按论文方法部分复核。
\textbf{证据与边界：} 本轮召回分类尚未逐项核对实验表格；后续证据补全阶段需要读取论文源码/PDF或项目页确认数据集、指标和对比方法。 该条来自宽召回批次；源码下载受 arXiv e-print 网络限制，部分论文暂时只能依据标题、摘要和类别分类，需要在入可信书前做人工/LLM精读复核。 本书只采用本地 paper pack、数据库分类、源码/引用登记和论文摘要中可核验的信息；没有本地证据支撑的速度、排行榜或开源完整性不在正文中强化。\textbf{方法信息：} 公开链接：\url{https://github.com/latex3/latex3.github.io}，完整性按本地证据保守处理；引用信号：235 citations (8.1/month, domain-relative cutoff 2.8/month, semantic\_scholar)。
\textbf{设计动机与边界：} 相对原始 3DGS 以独立高斯和图像重建误差为主的设定，3D Gaussian as a New Era 把问题落在以下可实现的改动上：从摘要层面看，3D Gaussian as a New Era 围绕“把 3DGS 用于机器人、导航或具身智能”展开。若其正文或代码后续可读，应进一步核对它具体改动的是表示、初始化、渲染、优化、损失、动态建模、压缩或应用接口中的哪一环。其设计目的不是单独增加一个模块，而是用上述表示、监督或渲染约束消除原流程在目标任务中的具体失配。该工作被召回是因为论文文本包含 Gaussian/3D Gaussian 相关表述；当前分类先确定其任务域和与 3DGS 的相关性，后续证据流水线再补代码、企业和引用信息。适用时还应保留论文给出的边界：该条来自宽召回批次；源码下载受 arXiv e-print 网络限制，部分论文暂时只能依据标题、摘要和类别分类，需要在入可信书前做人工/LLM精读复核。
% METHOD: 2501.01895
\paragraph{EnerVerse（4D 世界生成与动作预测）（2025.01）}
EnerVerse~\cite{2501_01895} 把视频生成、4DGS 和动作策略连接成具身未来空间模型。它的生成头采用分块自回归视频扩散，逐块预测未来空间；稀疏上下文记忆保留关键先验，支持长序列生成。

数据引擎将生成建模与 4DGS 结合，通过深度变形获得 3D 重建并渲染多视角图像，形成减少 sim-to-real 差距的数据循环。动作策略头复用去噪特征预测动作块，把 4D 世界表示转成物理动作。相对 GWM 的高斯潜空间未来预测，EnerVerse 更大规模地把视频扩散和 4DGS 用作机器人基础模型数据飞轮。
\textbf{实现细节：} EnerVerse（2025.01） 采用多视角视频表示结合 4DGS 数据引擎形成自强化数据循环——通过 4DGS 从机器人交互数据中重建动态场景，生成的新视角数据反过来用于训练更强大的视频生成模型，形成数据质量和模型能力的持续提升闭环。
\textbf{实现补充：} 实现上还要核对输入、输出、参数、特征、初始化、训练损失、正则、渲染或推理接口，以及优化中的剪枝、分裂、采样、编码或解码步骤；这些抓手决定该方法到底改了 3DGS 管线的哪一环。
\textbf{方法信息：} 引用很多，Semantic Scholar 约 58 次。
\textbf{设计动机与边界：} 相对原始 3DGS 以独立高斯和图像重建误差为主的设定，EnerVerse 把问题落在以下可实现的改动上：EnerVerse包含三个核心组件：1）EnerVerse-G（生成头）：采用分块自回归视频扩散框架，定义「块」为未来空间最小单元，迭代预测下一个块扩展具身空间。设计稀疏上下文记忆机制在训练中保留关键先验内容，防止模型崩溃并支持无限长度序列生成。使用多视角视频表示解决运动歧义和3D定位问题；2）EnerVerse-D（数据引擎）：结合生成建模与4DGS形成自强化数据循环，通过深度变形获取3D重建并渲染多视角图像，减少sim-to-real差距；3）EnerVerse-A（动作策略头）：将4D世界表示转化为物理动作，复用第一步去噪特征预测动作块，在单RTX 4090上实现约280ms/8步动作块的效率。其设计目的不是单独增加一个模块，而是用上述表示、监督或渲染约束消除原流程在目标任务中的具体失配。提出将视频生成与4DGS结合的具身机器人基础模型，通过分块自回归视频扩散预测未来具身空间，设计稀疏上下文记忆实现长期推理，结合4DGS形成数据飞轮减少sim-to-real差距，实现SOTA机器人操作性能。适用时还应保留论文给出的边界：多视角视频生成的计算开销大；数据飞轮需要4DGS重建质量保证；对未见过的任务泛化能力需进一步验证。
% METHOD: 2606.01458
\paragraph{LEGS（Embodied Gaussian Splatting World 中的 VLA 微调）（2026.05）}
LEGS~\cite{2606_01458} 旨在减少人形 loco-manipulation 对人工遥操作示范的依赖。它构建 Embodied Gaussian Splatting 混合仿真环境：真实部署场景由手持采集重建为 3DGS 背景，机器人、物体和道具作为 mesh 前景合成，再通过程序化动作原语生成交互数据。

为了让合成图像接近真实机器人相机，系统采用两阶段颜色校准对齐前景和背景外观。随后用这些合成交互对 VLA 策略微调，使策略学习在目标环境中导航和操作。相对 GaussGym 的大规模运动仿真，LEGS 更聚焦人形机器人和无需 teleoperation 的本地环境微调。

机制上，输入是目标房间或实验场地的短时采集视频；重建模块输出可渲染的 3DGS 背景，物理模块只负责可交互前景，策略训练模块再把渲染图像、动作原语和任务成功信号组成监督。这样做把真实外观、可控交互和策略优化分开：背景不需要完整物理仿真，前景不需要照片级重建，但合成相机流仍能进入同一个 VLA 微调流程。
\textbf{方法信息：} 代码：\url{https://github.com/AaronNZH/LeGS}。
\textbf{设计动机与边界：} 相对原始 3DGS 以独立高斯和图像重建误差为主的设定，LEGS 把问题落在以下可实现的改动上：系统从手持采集重建部署环境3DGS，将机器人、物体和道具作为mesh前景合成，并通过两阶段颜色校准对齐真实机器人相机。其设计目的不是单独增加一个模块，而是用上述表示、监督或渲染约束消除原流程在目标任务中的具体失配。提出Embodied Gaussian Splatting混合仿真环境，把真实场景3DGS背景、mesh前景和程序化动作原语结合起来，用于VLA策略微调和sim-to-real。适用时还应保留论文给出的边界：物理交互主要由mesh前景承担，3DGS背景不可交互；任务类型和场景范围有限。
% METHOD: 2606.03682
\paragraph{GN0（视觉语言导航的生成、评测和策略学习统一范式）（2026.06）}
GN0~\cite{2606_03682} 把 3DGS 用作视觉语言导航数据引擎。它构建从 3D 场景到导航数据、BEV 评测和策略学习的统一管线，其中 3DGS 提供照片级漫游和碰撞感知仿真环境。

系统生成 GN-Matrix 数据集，支持交互式 roaming 和动态 avatar 评测，再训练 Break and Establish 导航基础模型。相对 LagMemo/VISTA 等单系统导航方法，GN0 更像统一基准和数据生成平台：3DGS 的作用是让导航数据既可渲染、可交互，又有可评估的几何环境。
\textbf{实现细节：} GS-Verse（2025.10） 直接将物体网格与高斯泼溅表示集成，通过网格-高斯映射 $\Phi$ 建立高斯基元与网格顶点之间的关联，在 VR 环境中实现更精确的表面近似和高度真实的变形交互。18 人用户研究证实该方法在物理感知拉伸操控上统计显著优于 VR-GS 基线，同时在系统可用性（SUS）和认知负荷（NASA-TLX）方面表现优异。VR-Doh（2024.12） 定制物质点方法（MPM）模拟手诱导大变形，通过中轴变换（MAT）近似手部几何实现高效碰撞处理，结合解耦外观与物理表示和均匀高斯体积正则化技术增强 3DGS 无缝渲染，支持新手用户直观的接触式和手势式 3D 建模体验。
\textbf{实现补充：} 实现上还要核对输入数据、输出指标、基线配置、评测协议、系统接口和适用边界；若调用 3DGS 渲染或重建结果，也应说明相机参数、训练设置、优化流程和指标计算如何对齐。
\textbf{方法信息：} Meta 参与了这篇论文。
\textbf{设计动机与边界：} 相对原始 3DGS 以独立高斯和图像重建误差为主的设定，GN0 把问题落在以下可实现的改动上：系统构建GN-Matrix数据集，使用3DGS引擎支持交互式roaming和动态avatar评测，再训练Break and Establish导航基础模型。其设计目的不是单独增加一个模块，而是用上述表示、监督或渲染约束消除原流程在目标任务中的具体失配。提出从3D场景到导航数据、BEV评测和策略学习的统一管线，其中3DGS提供照片级漫游和碰撞感知仿真环境。适用时还应保留论文给出的边界：3DGS主要作为仿真平台底座；导航策略本身不是3DGS方法创新。
% METHOD: 2606.04158
\paragraph{Multi-Agent NBV（风险规避多机器人视角规划）（2026.06）}
Multi-Agent Next-Best-View Optimization for Risk-Averse Planning~\cite{2606_04158} 把 3DGS 用于多机器人主动建图。每个机器人维护局部 3DGS 地图，沿候选轨迹估计期望信息增益，但团队不共享原始传感器数据，只交换候选视角、轨迹描述和标量信息增益贡献。

分布式目标通过通信图上的 Consensus ADMM 求解，减少通信负担。路径风险由平均风险值建模，特别关注碰撞风险的尾部事件，而不是只优化平均距离或平均碰撞概率。相对 AIR-Embodied 的单体交互式重建，它关注多智能体协作和风险规避规划。
\textbf{实现细节：} Multi-Agent Next-Best-View Optimization for Risk-Averse Planning（2026.06） 延续 3DGS map 在"主动建图与NBV规划"上的基本设定，主要改变是：分布式risk-aware NBV框架，每个机器人维护私有局部3DGS地图，通过C-ADMM协同最大化轨迹掩码区域的信息增益。3DGS 的几何-渲染统一表示保持为底层地图。其设计动机是：提出通信高效的多智能体NBV优化框架，只交换候选视角、轨迹描述和标量EIG贡献，同时用AVaR建模路径碰撞风险。具体而言，每个机器人维护局部3DGS地图并估计沿候选轨迹的期望信息增益，团队通过通信图上的Consensus ADMM求解分布式目标，避免共享原始传感器数据。适用边界是：依赖局部3DGS地图质量和候选视角集合；通信拓扑变化和真实机器人延迟可能影响收敛。
\textbf{实现补充：} 实现上还要核对输入、输出、参数、特征、初始化、训练损失、正则、渲染或推理接口，以及优化中的剪枝、分裂、采样、编码或解码步骤；这些抓手决定该方法到底改了 3DGS 管线的哪一环。
\textbf{方法信息：} Amazon 参与了这篇论文。
\textbf{设计动机与边界：} 相对原始 3DGS 以独立高斯和图像重建误差为主的设定，Multi-Agent Next-Best-View Optimization for Risk-Averse Planning 把问题落在以下可实现的改动上：每个机器人维护局部3DGS地图并估计沿候选轨迹的期望信息增益，团队通过通信图上的Consensus ADMM求解分布式目标，避免共享原始传感器数据。其设计目的不是单独增加一个模块，而是用上述表示、监督或渲染约束消除原流程在目标任务中的具体失配。提出通信高效的多智能体NBV优化框架，只交换候选视角、轨迹描述和标量EIG贡献，同时用AVaR建模路径碰撞风险。适用时还应保留论文给出的边界：依赖局部3DGS地图质量和候选视角集合；通信拓扑变化和真实机器人延迟可能影响收敛。
% METHOD: 2606.24628
\paragraph{ArtiTwinSplat（可交互铰接物体数字孪生）（2026.06）}
ArtiTwinSplat~\cite{2606_24628} 从 RGB-D 视频自动恢复可交互铰接物体数字孪生。普通 3DGS 能重建几何和外观，但无法告诉机器人抽屉、门或机械部件如何运动；ArtiTwinSplat 的改变是从观测到的运动中无监督发现 articulation。

它用 3D Gaussian Splatting 保持物体几何和光度细节，同时跟踪视频中的部件运动，恢复部件结构和关节参数，并联合优化外观、几何和运动。输出不需要 CAD、仿真资产或人工标注，却可用于后续交互式操作。相对 GSWorld 的资产库，它更专注单个铰接物体的可交互恢复。
\textbf{实现细节：} ArtiTwinSplat（2026.06） 延续 3DGS 在"Interactive digital twins for manipulation"上的基本设定，主要改变是：3DGS 外观重建结合无监督 articulation 发现和运动学估计。3DGS 的几何-渲染统一表示保持为底层地图。其设计动机是：无需 CAD、仿真资产或人工标注，即可从 RGB-D 视频自动生成具备几何、外观和关节可交互性的物体模型。具体而言，以 3D Gaussian Splatting 保持几何和光度细节，再通过无监督 articulation discovery 从观察到的运动中恢复部件结构和关节参数，并联合跟踪与优化。适用边界是：依赖 RGB-D 观测和可见运动，对遮挡严重或未充分运动的部件可能受限；具体速度和失败模式需看全文。
\textbf{方法信息：} Microsoft 参与了这篇论文。
\textbf{设计动机与边界：} 相对原始 3DGS 以独立高斯和图像重建误差为主的设定，ArtiTwinSplat 把问题落在以下可实现的改动上：以 3D Gaussian Splatting 保持几何和光度细节，再通过无监督 articulation discovery 从观察到的运动中恢复部件结构和关节参数，并联合跟踪与优化。其设计目的不是单独增加一个模块，而是用上述表示、监督或渲染约束消除原流程在目标任务中的具体失配。无需 CAD、仿真资产或人工标注，即可从 RGB-D 视频自动生成具备几何、外观和关节可交互性的物体模型。适用时还应保留论文给出的边界：依赖 RGB-D 观测和可见运动，对遮挡严重或未充分运动的部件可能受限；具体速度和失败模式需看全文。
% METHOD: 2606.30645
\paragraph{VLK（重建场景中的合成交互监督）（2026.06）}
VLK~\cite{2606_30645} 针对人形机器人缺少同步数据的问题：训练需要第一视角图像、语言命令和机器人可执行轨迹，但真实采集昂贵。方法用 3DGS 重建 metric-scale 室内环境，再利用 privileged scene information 合成导航和物体交互轨迹。

在生成数据时，系统知道场景几何、物体位置和可达区域，因此可以自动产生语言指令、目标轨迹和交互动作，并把机器人第一视角渲染为训练观察。相对 LEGS 的本地环境 VLA 微调，VLK 更偏大规模合成监督；相对 GN0，它关注人形 loco-manipulation 而非通用视觉语言导航评测。
\textbf{实现细节：} VLK（2026.06） 延续 3DGS 在"Embodied learning data from Gaussian scenes"上的基本设定，主要改变是：用 3DGS 重建室内环境并渲染 egocentric observations，配合任务语言和运动轨迹。3DGS 的几何-渲染统一表示保持为底层地图。其设计动机是：填补缺少大规模 egocentric images、language commands 和 robot-compatible trajectories 同步数据的问题。具体而言，通过 3DGS 重建 metric-scale indoor environments，用 privileged scene information 合成导航和物体交互轨迹，并渲染第一视角观察作为训练监督。适用边界是：3DGS 是数据生成环境而非机器人算法核心；sim-to-real 差距和场景交互物理真实性仍需验证。
\textbf{实现补充：} 实现上还要核对输入、输出、参数、特征、初始化、训练损失、正则、渲染或推理接口，以及优化中的剪枝、分裂、采样、编码或解码步骤；这些抓手决定该方法到底改了 3DGS 管线的哪一环。
\textbf{技术实现补充：} VLK。VLK（2026.06） 延续 3DGS 在"Embodied learning data from Gaussian scenes"上的基本设定，主要改变是：用 3DGS 重建室内环境并渲染 egocentric observations，配合任务语言和运动轨迹。3DGS 的几何-渲染统一表示保持为底层地图。其设计动机是：填补缺少大规模 egocentric images、language commands 和 robot-compatible trajectories 同步数据的问题。具体而言，通过 3DGS 重建 metric-scale indoor environments，用 privileged scene information 合成导航和物体交互轨迹，并渲染第一视角观察作为训练监督。适用边界是：3DGS 是数据生成环境而非机器人算法核心；sim-to-real 差距和场景交互物理真实性仍需验证。
\textbf{方法信息：} Amazon 参与了这篇论文。

\begin{takeaway}
机器人章节可以用“从看见到行动”的链条理解。Scalable Real2Sim、Splatting Physical Scenes、GaussTwin 和 GSWorld 把 3DGS 接到物理仿真和数字孪生；LagMemo、VISTA、GaussExplorer 和 NavGSim 把高斯场变成导航记忆、探索目标和仿真后端；3DAffordSplat、SparseGrasp、Splat-MOVER 和 ViserDex 把语义、可供性和策略迁移接入机械臂；GWM、Physically Embodied GS 和 EnerVerse 则把高斯提升为动作条件世界模型。对于机器人任务，3DGS 的价值不止是照片级渲染，而是把真实世界压缩成一个可查询、可更新、可仿真、可规划的中间层。
\end{takeaway}
\textbf{设计动机与边界：} 相对原始 3DGS 以独立高斯和图像重建误差为主的设定，VLK 把问题落在以下可实现的改动上：通过 3DGS 重建 metric-scale indoor environments，用 privileged scene information 合成导航和物体交互轨迹，并渲染第一视角观察作为训练监督。其设计目的不是单独增加一个模块，而是用上述表示、监督或渲染约束消除原流程在目标任务中的具体失配。填补缺少大规模 egocentric images、language commands 和 robot-compatible trajectories 同步数据的问题。适用时还应保留论文给出的边界：3DGS 是数据生成环境而非机器人算法核心；sim-to-real 差距和场景交互物理真实性仍需验证。
